\documentclass[review,3p,authoryear]{elsarticle}

% --- Packages ---
\usepackage{graphicx}
\usepackage{booktabs}
\usepackage{amsmath}
\usepackage{amssymb}
\usepackage{amsthm}
\usepackage{bm}
\usepackage{mathrsfs}
\usepackage{siunitx}
\usepackage{natbib}
\usepackage{subcaption}
\usepackage{multirow}
\usepackage{lineno}
\usepackage{booktabs}
\usepackage{algorithm}
\usepackage{algpseudocode}
\usepackage{url}
\usepackage{placeins}
% \modulolinenumbers[]

\algrenewcommand\algorithmicrequire{\textbf{Input:}}
\algrenewcommand\algorithmicensure{\textbf{Output:}}

\newtheorem{theorem}{Theorem}[section]
\newtheorem{proposition}[theorem]{Proposition}
\newtheorem{corollary}[theorem]{Corollary}
\newtheorem{lemma}[theorem]{Lemma}

\journal{Transportation Research Part C: Emerging Technologies}

\begin{document}

\begin{frontmatter}

\title{Pricing Power and the Price of Disruption under Orbital Handoff in Multi-Operator LEO Satellite Communication and Routing Networks}
\author[a]{Dung-Ying Lin}
\author[b]{Surendra Reddy Kancharla\corref{cor1}} 
\ead{surendra_reddy.kanchalra@tu-dresden.de}

\author[b]{S. Travis Waller}


\cortext[cor1]{Corresponding author}
\address[a]{Department of Industrial Engineering and Engineering Management, National Tsing Hua University (NTHU), Taiwan}
\address[b]{"Friedrich List" Faculty of Transport and Traffic Sciences, Technische Universität Dresden, Saxony, Germany, 01069}

\begin{abstract}
This paper examines how orbital handoff events reshape market competition and welfare in multi-operator LEO satellite packet routing networks. We develop a bi-level framework in which non-atomic users route packets under Wardrop equilibrium while competing operators set per packet access prices in a Nash pricing game. A key structural result is that orbital events partition the market into two regimes: a competitive regime with sufficient rival capacity and a captive flow regime in which one operator acquires monopoly power over residual demand. Once captive demand becomes positive, uncapped prices become unbounded and no finite equilibrium exists. The resulting welfare loss, measured by the Price of Disruption (PoD), ranges from 0.67\% to 5.72\% in our experiments and is concentrated within 30 to 120 second disruption windows. Three policy insights emerge. First, competition collapse is exactly predictable from public orbital data through the rival capacity margin $M_{n}^{k}$. Second, pre-activating a price cap before the regime transition imposes no cost in the competitive regime while reducing welfare loss by up to 12.8\% in the captive flow regime. Third, equilibrium prices diverge at rate ${\Theta} \left({\varepsilon}^{-1/2}\right)$ as the margin approaches zero, providing a market observable early warning signal independent of operator count. The framework establishes six theoretical results for general $N$ and arbitrary network topology, together with closed form results for the two-operator case, including a scalar equilibrium equation and a quasi-concavity sufficient condition. Scalability experiments further show that the full per window computational pipeline can be executed within the 15 to 60 second orbital window for constellations with up to five competing operators.
\end{abstract}
\begin{keyword}
LEO satellite packet routing \sep Wardrop equilibrium \sep bi-level pricing \sep Price of Disruption \sep congestion pricing
\end{keyword}
\end{frontmatter}

\section{Introduction and Background}

Satellite internet now carries real-time data traffic for millions of users through networks of
hundreds to thousands of satellites in Low Earth Orbit (LEO). At the physical layer, a user's
data is broken into packets that traverse a sequence of satellite links from source to
destination, with each link contributing both propagation delay and queuing delay that grows
with traffic load. When multiple competing operators offer routing paths, users implicitly
make a route-choice decision that collectively determines the flow pattern across the network.
This flow pattern is precisely the object studied by \cite{wardrop1952road} and the transportation
science literature on network equilibrium: a distribution of demand over paths such that no
individual user can reduce their generalized cost by unilaterally switching routes.

The connection between packet routing and transportation network flow assignment is not
merely analogical. \cite{darroch1972} showed that M/M/1 queuing networks satisfy
the conditions of \cite{beckmann1956studies}, making the Wardrop potential function directly
applicable to data networks. \cite{bertsekas1987} established the equivalence
formally for wireline data networks. In the present LEO satellite setting, each competing
operator's set of satellite links constitutes a transport corridor, access prices function as
congestion tolls, and the regulator's design problem is structurally identical to the congestion
pricing and toll-setting problems studied in transportation science \citet{labbe1998bilevel,yang2004multi,jahn2005system}. The feature that distinguishes LEO satellite packet
routing from terrestrial transportation is that link capacities are not fixed but change
discontinuously at intervals of 15 to 60 seconds, driven by orbital mechanics. As a satellite
passes overhead, the elevation angle of the ground-satellite link rises and then falls. When
the angle drops below approximately 20 to 25 \textdegree{}, the link becomes commercially
unusable and deactivates; when a new satellite rises above the threshold, a new link activates.
This sequence of events, fully predictable from publicly available Two-Line Element (TLE)
orbital data, periodically reshapes the set of feasible routing paths and the capacity available
on each path.

The central finding of this paper is that these orbital events can trigger a qualitative change in the market structure of the routing game: once rivals jointly lack the capacity to carry a shifted flow, price competition collapses and total user routing cost jumps discontinuously at the transition, a welfare loss we term the Price of Disruption. Because this capacity schedule is known in advance, the collapse can be anticipated and pre-empted rather than merely diagnosed after the fact.

These phenomena sit at the intersection of three streams of the transportation literature, which the present paper connects and extends. The first is Wardrop equilibrium and network flow assignment. \cite{wardrop1952road, beckmann1956studies, daganzo1977traffic, sheffi1985urban} established the foundational user-equilibrium framework. \cite{roughgarden2002bad, roughgarden2005selfish} characterized efficiency loss in non-atomic routing games. The present paper extends this literature to a setting where the network topology, and hence the feasible set of the lower-level equilibrium problem, changes endogenously over time according to a deterministic physical law.

The second stream is bi-level congestion pricing. The bi-level programming framework for toll-setting was established by \cite{labbe1998bilevel} and extended by \cite{dewez2008new}, \cite{ban2006general}, and others. In these models a regulator (upper level) sets tolls, taking as given the users' Wardrop response (lower level). More recent work solves the resulting nonconvex pricing optimization with black-box surrogate methods such as Bayesian optimization, which can suffer from instability of the Gaussian-process surrogate \citep{huo2023bayesian}; the regulatory cap we study instead reduces to a closed-form bound (Theorem~\ref{thm:qcsc}) or a one-dimensional search, avoiding high-dimensional surrogate optimization. In the present paper the upper level is occupied
not by a regulator but by competing operators, each maximizing its own revenue. The
regulator's role is secondary: to impose the minimum price cap that prevents Nash
breakdown. This is closer in spirit to the access pricing literature \citet{Armstrong2001, laffont1993theory} but applied to a physical routing network.

The third stream is queueing network equilibrium. The connection between M/M/1 queueing networks and Wardrop equilibrium was established by \cite{darroch1972} and formalized by \cite{bertsekas1987}. \cite{johari2004efficiency,acemoglu2007competition} studied efficiency loss in M/M/1 pricing games. We extend this stream by incorporating topology dynamics, in which the M/M/1 capacity vector changes discontinuously over time due to orbital handoff events. We further characterize how these topology changes induce regime transitions in the Nash pricing game and provide the corresponding bi-level pricing analysis.

Beyond these three streams, a largely separate body of work treats satellite and air-traffic resource allocation as deterministic scheduling. Satellite range and multi-satellite observation scheduling have been formulated as integer programs solved by exact methods and metaheuristics \citep{chen2019mixed, wang2025solving}, and nanosatellite task scheduling has been addressed by Dantzig--Wolfe branch-and-price to improve mission quality-of-service \citep{rigo2022branch}; the common primitive (assigning activities to fixed, non-overlapping availability windows) is the classical fixed-interval scheduling problem \citep{kovalyov2007fixed}, whose applications include bandwidth allocation on communication channels. Closely analogous time-windowed scheduling arises in air-traffic management, notably the terminal-airspace scheduling problem \citep{ng2024optimization}. This literature optimizes how a \emph{single} operator assigns its tasks or transmissions to known visibility windows; we instead ask how \emph{competing} operators price access and how users route over the same time-varying capacity in equilibrium, which makes it complementary to our analysis.

Against this background, this paper makes three sets of contributions to the transportation science literature on Wardrop equilibrium and congestion pricing. \emph{Theoretically}, we first establish that the standard Wardrop equilibrium framework from transportation network analysis applies without modification to multi-operator LEO satellite packet routing once M/M/1 delay functions replace BPR functions; the M/M/1 hard-capacity barrier is not a modeling convenience but a structural necessity, since it is what creates the captive-demand floor ${\chi_n}$ that drives every subsequent result. On this foundation we prove three results that hold for general $N$ and arbitrary network topology: the Competition Collapse theorem (Theorem~\ref{t2a}), which shows that the market ceases to admit a finite Nash equilibrium precisely when the rival-capacity margin $M_{n}^{k} < 0$, a condition computable from TLE data without solving any equilibrium problem; the Welfare Discontinuity theorem (Theorem~\ref{t3}), which shows that total user routing cost jumps upward at every A-to-B topology transition, so PoD > 0 at each disruption, the satellite-routing analogue of the price-of-anarchy results of \citet{roughgarden2002bad}, but driven by topology change rather than selfish routing per se; and an asymptotic price-divergence result (Proposition~\ref{p3}), which shows that equilibrium prices grow at rate ${\Theta}({\varepsilon}^{-1/2})$ as the margin approaches zero, giving a TLE-computable early-warning indicator that is independent of $N$ and requires only observed market prices rather than full equilibrium computation.

\emph{Methodologically}, for the two-operator case we reduce the Nash pricing problem to a scalar equilibrium equation solvable by one-dimensional root finding, together with a closed-form sufficient condition for quasi-concavity (Theorem~\ref{thm:qcsc}) expressed entirely in network topology parameters; these are the satellite-routing analogues of the bi-level toll-setting results of \cite{labbe1998bilevel} and \cite{dewez2008new}. \emph{For policy}, we derive the minimum regulatory price cap that restores Nash existence (Theorems~\ref{t2a} and~\ref{t2b}), show that it admits a closed form in topology parameters for $N$= 2 and an $O(1)$ per-window evaluation from TLE data, and establish that pre-activating it at the last competitive window before a predicted transition is costless in that regime (the cap does not bind while $M_{n}^{k} \ge 0$) yet prevents the collapse, making proactive capping the dominant regulatory strategy. We quantify these effects with the Price of Disruption (PoD) metric through $N$= 3, 4, 5 numerical experiments and a calibrated Taiwan gateway case study, in which orbital disruptions generate PoD up to 5.72 percent and proactive capping reduces peak-disruption cost by up to 12.8 percent relative to the uncapped case.

The general framework of Sections \ref{sec2} and \ref{sec3} covers $N$ competing operators and any network topology; Theorems \ref{t1} through \ref{t3} and Propositions \ref{p2} through \ref{p4} are the main theoretical contribution. Section \ref{sec4} specializes to the two-operator, two-path parallel network, the minimal model that permits closed-form expressions for all key quantities. Section \ref{sec5} presents the algorithms and Section \ref{sec6} the computational study, where Section \ref{sec6.7} confirms that Theorem \ref{t3} and Proposition \ref{p2} extend to a 20-node network with $N$= 4 operators.

\section{Mathematical Model}
\label{sec2}

This section formalizes the packet routing network as a Wardrop-Nash bi-level game. The
lower level is a standard Wardrop user equilibrium, the upper level a Nash pricing game
among competing operators. The key departure from standard bi-level congestion pricing
models is that the lower-level feasible set, the capacity vector ${\mu}_{k}$, changes over time according
to orbital mechanics, not at the regulator's discretion.

\subsection{The LEO Satellite Packet Routing Network}

For readers from the transportation science community, we briefly describe the physical network before formalizing it. Imagine a road network in which the lanes on each road segment open and close periodically according to a publicly known schedule; the satellite problem is its exact digital analogue.

The network $G(t) = (V, E(t))$ consists of ground stations (origins and destinations of packet flows) and satellites (intermediate relay nodes). A typical ground station generates and receives data at aggregate rates in the hundreds of Mbps to Gbps range. A gateway aggregates traffic from many end-users and connects to the satellite constellation. Two types of links exist: ground-satellite links (GSLs), active when elevation angle ${\varepsilon}_{\mathrm{gs}}(t) \ge {\varepsilon}_{\min}$ (typically 20--25\textdegree{}), and inter-satellite links (ISLs), active when satellite separation $d_{ss'}(t) \le d_{\max}$. A link's capacity is not a fixed parameter but a function of orbital geometry: ${\mu}({\varepsilon}) = {\mu}_{\max} \cdot \sin({\varepsilon})/\sin({\varepsilon}_{\mathrm{ref}})$ per ITU-R S.1528 \citet{ITUR}. As a satellite rises and sets, the capacity on its GSLs traces a smooth arc from zero to peak and back to zero, before a discrete deactivation at ${\varepsilon} = {\varepsilon}_{\min}$.

Data packets from many independent users traverse paths from ground stations through the constellation. Because each individual packet is small relative to total traffic, users are modeled as non-atomic: each user controls an infinitesimal fraction of total demand $D$. Users route packets to minimize individual generalized cost: the sum of propagation delay and queuing delay on each traversed link, plus access tolls charged by the operator controlling each link. The resulting flow pattern is a Wardrop user equilibrium.

Two-Line Element (TLE) sets, published by the U.S. Space Force at celestrak.org and updated multiple times daily, specify the orbital parameters of every tracked satellite. From TLE data, the full future trajectory of link activations, deactivations, and capacity arcs is deterministically computable hours to days in advance, a feature exploited directly in the algorithm designs of Section \ref{sec5} and the policy results of Section \ref{sec6}.

\subsection{Network Topology and Quasi-Static Windows}

The LEO satellite packet routing network is modeled as a time-varying directed graph $G(t) = (V, E(t))$ with vertex set $V$ partitioned into ground stations and satellites. There are $N$ competing operators indexed by $n {\in} \mathscr{N} = \{1, {\dots}, N\}$; operator $n$ has exclusive control of link set $E_{n}$. Operators' link sets are disjoint: $E_{n} {\cap} E_{m} = {\varnothing}$ for $n {\ne} m$. Total demand $D$ > 0 is the aggregate packet traffic from a fixed source to a fixed destination, held constant within each window (the analysis extends to multiple OD pairs by standard superposition).

We discretize time into quasi-static windows $T_{k}$ = [$t_{k}$, $t_{k+1}$) of 15 to 60 seconds. Within window $k$, the capacity vector ${\mu}_{k}$ = (${\mu}_{1,k}$, \dots, ${\mu}_{N,k}$) is fixed. At each window boundary a topology event updates ${\mu}_{k}$, reflecting link activations and deactivations driven by orbital mechanics.

The delay on link $e$ controlled by operator $n$, with capacity ${\mu}_{e}$ (packets per unit time) and flow $f_{e}$ (packets per unit time), is:
\begin{equation}
    l_e (f_e, \mu_e) = t_{e} + \frac{1}{\mu_{e} - f_{e}}, f_e \in [0, \mu_e)
    \label{eq1}
\end{equation}

The first term $t_{e}$ is propagation delay (determined by orbital geometry, not traffic). The second term is M/M/1 queuing delay, which grows without bound as flow approaches capacity. The M/M/1 specification is standard for satellite link-level queueing \citet{bertsekas1987} and is structurally necessary for this paper: the hard capacity barrier $f_{e} < {\mu}_{e}$ is what creates the captive demand floor ${\chi}_{n}$ in Regime B. BPR or other soft-capacity delay functions allow flows to exceed ${\mu}_{e}$, preventing the definition of captive demand and invalidating the competition-collapse mechanism.

\subsection{Wardrop User Equilibrium: The Packet-Routing Lower Level}

Each user (a packet or, in the continuum model, an infinitesimal fraction of traffic) selects a
routing path from source to destination to minimize individual generalized cost: the sum of
delays on all traversed links plus the per-packet access tolls charged by the operator
controlling each link. Because users are non-atomic, the flow pattern induced by individually
optimal routing satisfies the \cite{wardrop1952road} user-equilibrium condition: all used paths
between any OD pair carry equal generalized cost, and this cost does not exceed the cost on
any unused path.

Following \cite{beckmann1956studies}, the Wardrop user equilibrium is equivalently characterized as the minimizer of the Beckmann potential:
\begin{equation}
    \Phi(f,p; \mu_k) = \sum_{e \in E_k} t_e f_e + p_{n(e)} f_e - \mu_e \ln \left(1 - \frac{f_e}{\mu_e}\right)
    \label{eq:2}
\end{equation}

where the summation is over all active links, $p_{n(e)}$ is the per-packet toll charged by the operator controlling link $e$, and the logarithmic term is the anti-derivative of the M/M/1 queuing delay. Comparing equation \eqref{eq:2} with the standard Beckmann potential for road networks: the M/M/1 log-barrier term replaces the BPR integral. The formal structure is identical, a strictly convex program whose minimizer characterizes the Wardrop equilibrium, but the barrier creates the capacity constraint that is absent from Bureau of Public Roads (BPR) models.

For the two-operator, two-path specialization of Sections \ref{sec4} and \ref{sec5}, the flow space reduces to $(f_{1}, f_{2})$ subject to $f_{1} + f_{2} = D$, $0 < f_{1} < {\mu}_{1}$, $0 < f_{2} < {\mu}_{2}$. This two-path parallel network is the satellite-routing analogue of a two-route road network with congestion tolls, the canonical setting of \cite{Pigou1920} and many subsequent studies in transportation.

\subsection{The Operator Nash Pricing Game: The Upper Level}

Operator $n$ sets a per-packet access toll $p_{n} {\in} [0, \bar{p}]$, where \ensuremath{\bar{p}} is a regulatory price cap (\ensuremath{\bar{p}} = \ensuremath{\infty} if no regulation), to maximize revenue ${\pi}_{n}(p) = p_n {\cdot} F_n*(p; {\mu}_k)$, where $F_{n}^{*}$ is the total flow through operator $n$'s links at the Wardrop equilibrium induced by price vector $p$. Operators act simultaneously and non-cooperatively, taking rivals' prices as fixed (Nash conjecture).

The rivalry between operators is mediated through the Wardrop routing of users. When operator $n$ raises its price, users route more packets through rivals' links, reducing $F_{n}^{*}$. This is the standard demand-substitution mechanism of Bertrand competition, transplanted to a packet-routing network. The critical new element is that the degree of substitutability is determined by orbital geometry: when rivals collectively lack capacity to carry all packets ($\ensuremath{\bar{{\mu}}}_{-n}^{k}$ < $D$), substitution becomes impossible for a fraction ${\chi}_{n}$ of demand.

To characterize this formally, define the rival capacity of operator $n$ in window $k$ as the maximum packet flow that all rival operators can jointly carry, excluding $n$'s links:
\begin{equation}
    \bar{{\mu}}_{-n}^{k} := \max\left\{ \text{total } s\text{--}d \text{ flow} \;:\; f_e \le {\mu}_e \;\; \forall e \in \bigcup_{m \ne n} E_m, \;\; f_e = 0 \;\; \forall e \in E_n \right\}
\end{equation}

This is a standard max-flow LP. For the $N$-operator parallel network it reduces to $\ensuremath{\bar{{\mu}}}_{-n}^{k}$ = $\sum_{m {\ne} n} {\mu}_{m,k}$, requiring no LP. The rival capacity margin and captive demand floor are:
\begin{equation}
    M_{n}^{k} := \frac{\ensuremath{\bar{{\mu}}}_{-n}^{k}}{D} - 1, \quad {\chi}_{n} := \max\{0, D - \ensuremath{\bar{{\mu}}}_{-n}^{k}\}
\end{equation}

These two quantities partition the orbital window into two market regimes, which are the central objects of analysis throughout the paper:

\textbf{Regime A (Competitive):} $M_{n}^{k}$ \ensuremath{\ge} 0 for all $n$; ${\chi}_{n}$ = 0. Rivals have sufficient joint capacity to carry all packets without operator $n$; every packet is contestable.

\textbf{Regime B (Captive Flow):} $M_{n}^{k} < 0$ for some $n$; ${\chi}_{n}$ = $D$ - $\ensuremath{\bar{{\mu}}}_{-n}^{k} > 0$. Rivals jointly lack capacity to carry all packets; operator $n$ holds a monopoly over a fraction ${\chi}_{n}$/$D$ of total packet demand regardless of its price.

\section{General Theoretical Results}
\label{sec3}

This section establishes six results that hold for any $N$ competing operators and any network topology. The transportation reader will recognize the proof strategy: Theorem \ref{t1} and Theorems \ref{t2a}--\ref{t2b} invoke the same fixed-point and convexity arguments used for road network equilibrium existence; Theorem \ref{t3} uses the continuity established in Theorem \ref{t1}. The new content is the analysis of how the orbital topology event (the change in ${\mu}_{k}$) moves the system between regimes and what this implies for welfare.

\subsection{Wardrop Equilibrium Existence and Continuity}

Before analyzing the pricing game, we must verify that the lower-level Wardrop equilibrium
is well-defined at every topology window and price vector. In road-network models this is guaranteed by strict convexity of the Beckmann potential \citet{beckmann1956studies, daganzo1977traffic}. The following theorem confirms the same property for the M/M/1 satellite routing potential, extending the result to $N$ operators.

The continuity results (parts (b) and (c)) are less standard but critical for the analysis that follows: they guarantee that the upper-level revenue functions are well-defined and continuous, enabling the existence of Nash equilibrium in part (d). In the road-network pricing literature, analogous continuity results are typically assumed; here they follow from the strict convexity of the M/M/1 potential and Berge's Maximum Theorem.

\begin{theorem}[Wardrop Equilibrium for $\bm{N}$ Operators]\label{t1}
Fix a quasi-static window $k$. Let $p{\in}[0,\bar{p}]^{N}$ be any price vector and let ${\mu}_{k}$ be any capacity vector such that the lower-level feasible set $\mathscr{F}({\mu}_{k})$ is nonempty. Then:

(a) A unique Wardrop user equilibrium $f^{*}(p;{\mu}_{k})$ exists.

(b) The equilibrium flow mapping $p {\mapsto} f^{*}(p;{\mu}_{k})$ is continuous on $[0, \bar{p}]^{N}$.

(c) For each operator $n$, the revenue function ${\pi}_{n}(p)=p_{n}F_{n}^{*}(p;{\mu}_{k})$ is continuous in $p$.

(d) Under Assumption QC, namely that ${\pi}_{n}(p_{n},p_{-n})$ is quasi concave in $p_{n}$ for every fixed $p_{-n}$, the upper level pricing game admits a pure strategy Nash equilibrium $p^{*}{\in}[0,\bar{p}]^{N}$.
\end{theorem}

\medskip\noindent\textit{Proof.} Consider the Beckmann potential

\begin{equation}
\tilde{{\Phi}}(f;p,{\mu}_{k})=\sum_{e{\in}E_{k}}^{} \left[t_{e}f_{e}-{\mu}_{e}\ln\left(1 - \frac{f_{e}}{{\mu}_{e}}\right)\right]+\sum_{n=1}^{N} p_{n}F_{n}(f),
\end{equation}

defined on the feasible set $\mathscr{F}({\mu}_{k})$. The linear toll term does not affect convexity. For each link $e$,

\begin{equation}
\frac{{\partial}^{2}\tilde{{\Phi}}}{{\partial}f_{e}^{2}}=({\mu}_{e}-f_{e})^{-2}>0 \quad\text{for all } 0\le f_{e}<{\mu}_{e}.
\end{equation}

Hence $\tilde{{\Phi}}({\cdot};p,{\mu}_{k})$ is strictly convex in $f$ on $\mathscr{F}({\mu}_{k})$. Since $\mathscr{F}({\mu}_{k})$ is nonempty and the logarithmic barrier satisfies

\begin{equation}
-{\mu}_{e}\ln\left(1 - \frac{f_{e}}{{\mu}_{e}}\right)\to+\infty \quad\text{as } f_{e}\uparrow{\mu}_{e}
\end{equation}

which implies that any minimizing sequence remains strictly within the feasible region $f_{e}<{\mu}_{e}$. Any minimizing sequence is bounded away from the capacity boundary. Therefore a minimizer exists, and strict convexity implies that it is unique. This proves part (a).

For part (b), the feasible set $\mathscr{F}({\mu}_{k})$ does not depend on $p$, and the objective $\tilde{{\Phi}}(f;p,{\mu}_{k})$ is jointly continuous in $\left(f, p\right)$. By part (a), the minimizer is unique for every $p$. Hence, by Berge's Maximum Theorem, the solution mapping $p{\mapsto}f^{*}(p;{\mu}_{k})$ is continuous.

Part (c) follows immediately, because $F_{n}^{*}(p;{\mu}_{k})$ is a continuous function of $f^{*}(p;{\mu}_{k})$, and ${\pi}_{n}(p)=p_{n}F_{n}^{*}(p;{\mu}_{k})$ is the product of continuous functions.

For part (d), each operator's strategy set $\left[0, \bar{p}\right]$ is compact and convex, so the joint strategy space $[0, \bar{p}]^{N}$ is also compact and convex. By part (c), each payoff ${\pi}_{n}(p_{n},p_{-n})$ is continuous in $p$. Under Assumption QC, ${\pi}_{n}$ is quasi concave in own strategy $p_{n}$, so each best response correspondence is nonempty, convex valued, and upper hemicontinuous. Kakutani's fixed point theorem therefore implies the existence of a pure strategy Nash equilibrium $p^{*}{\in}[0,\bar{p}]^{N}$. $\qed$

Theorem~\ref{t1} provides the foundation for everything that follows. The most important implication is that the orbital handoff problem is not one of equilibrium non-existence per se: at every ${\mu}_{k}$ with sufficient total capacity and any price cap, a Wardrop equilibrium and (under QC) a Nash pricing equilibrium exist. The question is whether the Nash equilibrium is competitive, which requires $M_{n}^{k}$ \ensuremath{\ge} 0 (Regime A), or whether it breaks down without a cap because competition collapses (Regime B). The next two theorems address this.

\subsection{Competition Collapse and Restoration}

The following pair of theorems constitutes the paper's central contribution. Theorem \ref{t2a} identifies the exact structural condition under which competition collapses: the rival capacity margin $M_{n}^{k} <$ 0. This is not a parametric condition (e.g., a threshold on demand or cost) but a topological one, determined purely by the orbital geometry at the current window. Theorem \ref{t2b} identifies the minimum regulatory intervention that restores competition.

In transportation science terms, Theorem \ref{t2a} is an existence-of-dominant-strategy result. When ${\chi}_{n} >$ 0, operator $n$'s routing revenue from captive packets alone grows without bound in its price. Raising price is always profitable: the captive packets cannot go elsewhere. This is analogous to a toll road with no untolled alternative, except that in the satellite case the condition is time-varying and its onset is predictable.

\begin{theorem}[Competition Collapse]\label{t2a}
Suppose $\ensuremath{\bar{{\mu}}}_{-n}^{k}$ < $D$ (Regime B), so ${\chi}_{n}$ = $D$ - $\ensuremath{\bar{{\mu}}}_{-n}^{k}$ > 0. For any fixed rival prices $p_{-n} {\ge} 0$:

\begin{equation}
{\pi}_{n}(p_n, p_{-n}) \ge p_n {\cdot} {\chi}_n \to +\infty \quad\text{as } p_n \to +\infty
\end{equation}

No finite best response exists for operator $n$. No finite pure-strategy Nash equilibrium exists (with or without a cap, unless the cap is finite).
\end{theorem}

\medskip\noindent\textit{Proof.} By the max-flow min-cut theorem applied to the rival subgraph (all links except $E_{n}$), the maximum packet flow achievable by all rivals jointly is exactly $\ensuremath{\bar{{\mu}}}_{-n}^{k}$. For any feasible flow, conservation at the destination node requires $F_{n}^{*}(p) + F_{-n}^{*}(p) = D$. Therefore $F_{n}^{*}(p) = D - F_{-n}^{*}(p) {\ge} D - \ensuremath{\bar{{\mu}}}_{-n}^{k} = {\chi}_{n} > 0$ for all finite $p_{n}$. The captive packets have nowhere else to go; they must route through $n$'s links. It follows that ${\pi}_{n} = p_{n} F_{n}^{*} \ge p_{n} {\chi}_{n} \to +\infty$ as $p_{n} \to +\infty$. $\qed$

Theorem \ref{t2a} shows that a regulator who waits until competition collapses before imposing a cap has already missed the opportunity: no finite Nash equilibrium exists, and packet routing costs are effectively unbounded. Theorem \ref{t2b} shows how a proactively designed cap restores the competitive equilibrium.

\begin{theorem}[Restoration of Nash Equilibrium under Regulatory Price Cap]\label{t2b}
Fix a quasi-static window $k$. Suppose the strategy space is restricted by a finite regulatory price cap, namely $p_{n}{\in}[0,\bar{p}],{\forall}n=1,{\dots},N,$ with $\bar{p}<{\infty}$. If, for each operator $n$, the revenue function ${\pi}_{n}(p_{n},p_{-n})$ is continuous in $p$ and quasi concave in its own strategy $p_{n}$ on $\left[0, \bar{p}\right]$, then the upper-level pricing game admits at least one pure strategy Nash equilibrium $p^{*}{\in}[0,\bar{p}]^{N}.$
\end{theorem}

A sufficient condition ensuring quasi concavity is provided by Theorem~\ref{thm:qcsc} in Appendix C for the two-operator case.

\subsection{Welfare Discontinuity and the Price of Disruption}

The third main result concerns what happens to total user routing cost at the moment of the A-to-B transition. In congestion pricing models, welfare is typically continuous in demand and in toll levels; discontinuities arise only in degenerate cases. Here, the topology event is itself a discrete shock, and the following theorem shows that it produces a discontinuous upward jump in total routing cost regardless of the pricing behavior of operators. Define total user routing cost at Nash equilibrium $p^{*}$(${\mu}_{k}$) as the sum of generalized costs paid by all packets:
\begin{equation}
TC({\mu}_{k}) = \sum_{e \in E_k} f_e^* \cdot l_e (f_e^*, \mu_k)
\end{equation}

and the Price of Disruption (PoD) at the topology event $k^*$ as the percentage jump:
\begin{equation}
    PoD(k^*) := \frac{TC({\mu}_{k^*}) - TC({\mu}_{k^*-1})}{TC({\mu}_{k^*-1})}, \quad k^* = \operatorname*{arg\,min}\{k : M_{n}^{k} < 0\}
    \label{eq6}
\end{equation}

PoD is closely related to the Price of Anarchy \cite{roughgarden2002bad}, but measures the welfare impact of a structural market change (topology transition) rather than the welfare impact of selfish routing per se. Positive PoD arises not because users route selfishly, they do that in both regimes, but because the topology change destroys competitive substitutability.

\begin{theorem}[Welfare Discontinuity at the Regime Transition]\label{t3}
Fix an operator $n$ and let $k^{*}$ denote a topology transition window such that $\bar{{\mu}}_{-n}^{k^{*}-1}=D+{\varepsilon},{\varepsilon}>0$ in Regime A, and $\bar{{\mu}}_{-n}^{k^{*}}<D$ in Regime B.

Assume a finite regulatory price cap $\bar{p}<{\infty}$. Then the total equilibrium user cost exhibits a strict upward jump at the transition:

\begin{equation}
\lim_{{\varepsilon}{\downarrow}0^{+}}TC^{*} \left({\mu}_{k^{*}-1}({\varepsilon})\right)<TC^{*}({\mu}_{k^{*}})
\end{equation}

Consequently, the Price of Disruption satisfies $PoD>0$ at every A to B topology transition, for any number of operators $N$.
\end{theorem}

\medskip\noindent\textit{Proof.} Approaching the transition from Regime A, let $\bar{{\mu}}_{-n}^{k^{*}-1}=D+{\varepsilon},{\varepsilon}{\downarrow}0^{+}$. Then by Proposition \ref{p3}, the equilibrium price satisfies $p_{n}^{*}({\varepsilon})={\Theta} \left({\varepsilon}^{-1/2}\right)$ and diverges as ${\varepsilon}{\to}0^{+}$. At the transition window $k^{*}$, we have $\bar{{\mu}}_{-n}^{k^{*}}<D$ so the captive demand is strictly positive:

\begin{equation}
{\chi}_{n}=D-\bar{{\mu}}_{-n}^{ k^{*}}>0
\end{equation}

Hence at least ${\chi}_{n}$ units of flow must be routed through operator $n$'s links under any feasible equilibrium flow. Under any finite cap $\bar{p}$, the equilibrium exists by Theorem \ref{t2b}, and the total equilibrium cost must satisfy

\begin{equation}
TC^{*}({\mu}_{k^{*}}){\ge}TC_{\min}+{\chi}_{n}(p_{n}^{*}+\min_{e{\in}E_{n}}t_{e})
\end{equation}

where $TC_{\min}$ denotes the minimum feasible delay cost absent the captive flow surcharge.

Because ${\chi}_{n}>0$ and $p_{n}^{*}{\ge}0$, the additional term is strictly positive, which yields

\begin{equation}
TC^{*}({\mu}_{k^{*}})>\lim_{{\varepsilon}{\downarrow}0^{+}}TC^{*} \left({\mu}_{k^{*}-1}({\varepsilon})\right)
\end{equation}

Therefore, the inequality is strict, and $PoD>0$. $\qed$

Theorem \ref{t3} has an important corollary for network design: in a multi-operator LEO satellite routing network, there is no price mechanism that makes welfare continuous across orbital topology events. The discontinuity is structural, not a consequence of regulatory failure. What regulation can do is minimize the magnitude of the jump: the theorem-based cap of Theorem \ref{t2b} limits the exploitation of captive demand, reducing PoD from its uncapped level. Sections \ref{sec6.2} through \ref{sec6.6} quantify this reduction numerically.

\subsection{Regime Transition Criterion}

Theorem \ref{t3} establishes that every A-to-B transition produces positive PoD. Proposition \ref{p2} provides the TLE-computable criterion that predicts exactly when such a transition occurs. In road-network terms, this is analogous to a capacity-drop event \citet{hall1987interpretation} that is fully predictable from physics, enabling pre-emptive traffic management.

\begin{proposition}[TLE Computable Regime Transition Criterion]\label{p2}
Fix an operator $n$. A topology event occurring at window $k^{*}$ triggers a transition from Regime A to Regime B for operator $n$ if and only if

$M_{n}^{k^{*}-1}{\ge}0$ and $M_{n}^{k^{*}}<0$

where $M_{n}^{k}=\frac{\bar{{\mu}}_{-n}^{k}}{D}-1$ denotes the rival capacity margin in window $k$. Equivalently, the transition occurs if and only if

$\bar{{\mu}}_{-n}^{k^{*}-1}{\ge}D$ and $\bar{{\mu}}_{-n}^{k^{*}}<D$

That is, the aggregate rival capacity crosses the demand threshold at $k^{*}$.

Because $\bar{{\mu}}_{-n}^{k}$ is fully computable from public Two-Line Element (TLE) orbital data for every future window $k$, the transition time $t^{k^{*}}$ is predictable in advance. For $N$-operator parallel networks,

\begin{equation}
M_{n}^{k}=\frac{\sum_{m{\ne}n}^{} {\mu}_{m,k}}{D}-1
\end{equation}

which is computable in $O(1)$ time per window.
\end{proposition}

Proposition \ref{p2} transforms the regulatory problem from reactive crisis management to proactive scheduling. Given the TLE data for the next 24 hours, one can compute the full schedule of A-to-B and B-to-A transitions and pre-activate the appropriate cap at each window. The next section quantifies the welfare benefit of this approach.

\subsection{Additional Results}

Two further results complete the theoretical picture; their proofs are in Appendix B.

\begin{proposition}[Asymptotic Price Divergence]\label{p3}
For any $N$-operator parallel network, if $\bar{{\mu}}_{-n}=D+{\varepsilon},{\varepsilon}{\downarrow}0^{+},$ then any interior profit maximizing equilibrium price satisfies

\begin{equation}
p_{n}^{*}({\varepsilon})={\Theta} \left({\varepsilon}^{-1/2}\right).
\end{equation}

Moreover, the divergence rate is independent of $N$. (Proof: Appendix B.1)
\end{proposition}

\begin{proposition}[N-Operator Equilibrium System]\label{p4}
For an $N$-operator parallel network in Regime A, the Nash equilibrium reduces to an $\left(N - 1\right)$-dimensional equilibrium system $\mathscr{E}_{N}$ in the flow ratios

\begin{equation}
r_{n}=\frac{f_{n}}{f_{1}},n=2,{\dots},N.
\end{equation}

System $\mathscr{E}_{N}$ admits a unique solution for any $N$. For $N=2$, it reduces to a scalar equilibrium equation. (Full statement and proof: Appendix B.2)
\end{proposition}

Proposition \ref{p3} serves as an early-warning indicator: equilibrium prices in the last few Regime A windows before the tipping point accelerates at a quantifiable rate, detectable from market data alone (see Section \ref{sec8.1}). Proposition \ref{p4} establishes that Algorithm \ref{alg3} converges to the unique Nash equilibrium for any $N$, grounding the computational study of Section \ref{sec6}.

\subsection{Two-Operator Specialization}
\label{sec4}

For $N=2$, the equilibrium system $\mathscr{E}_{N}$ of Proposition \ref{p4} reduces to a single scalar equilibrium equation in the flow ratio $r=f_{2}/f_{1}$ (Equation \eqref{eq9}, Appendix C), the only case admitting a one-dimensional reduction. By Corollary \ref{c3} this equation has a unique root $r^{*}$, from which the full Nash equilibrium is recovered in $O(\log\frac{1}{{\varepsilon}})$ time via Brent's method (Algorithm \ref{alg5}). The simplification also yields a closed-form minimum price cap $\bar{p}^{\mathrm{QC}}$ that guarantees Nash existence in Regime A, an $O(1)$-computable function of the topology parameters (Theorem~\ref{thm:qcsc}, Appendix C). The welfare-optimal regulatory cap solves
\begin{equation}
    \min_{\bar{p} \ge 0} TC(\bar{p}) \quad\text{s.t.}\quad \bar{p} \le \bar{p}^{\mathrm{QC}}(\mu_k),
    \label{eq7}
\end{equation}
computed numerically by Algorithm \ref{alg4}. For $N>2$, Nash equilibrium prices are computed by Algorithm \ref{alg3} (best-response iteration) and the QC condition is verified numerically per window.

\section{Algorithm Design}
\label{sec5}

We implement five algorithms that together constitute the regulatory toolkit. The algorithms are designed to be executable per orbital window from TLE input, with computation times well within the 15 to 60 second window length. Algorithms \ref{alg1} through \ref{alg3} are applicable for general $N$; Algorithms \ref{alg4} and \ref{alg5} exploit the closed-form structure available for $N$= 2. Input to the system per window: capacity vector ${\mu}_{k}$, propagation delay vector $t_{k}$, total demand $D$, price cap \ensuremath{\bar{p}}. Output: Nash equilibrium ($p^{*}$, $f^{*}$), regime classification (A or B), welfare metrics $\mathrm{TC}$ and $\mathrm{PoD}$.

\begin{algorithm}[!h]
\caption{Computation of the Wardrop Equilibrium for General $N$}
\label{alg1}
\begin{algorithmic}[1]
\Require Price vector $p$, capacity vector ${\mu}$, propagation delay vector $t$, total demand $D$
\Ensure Unique Wardrop equilibrium flow vector $f^{*}(p;{\mu})$
\State Initialize $f^{\left(0\right)}$ proportional to ${\mu}$, scaled such that $\sum_{n=1}^{N} f_{n}^{\left(0\right)} = D$.
\State Solve the strictly convex optimization problem
\Statex $\min_{f} \tilde{{\Phi}}(f,p)=\sum_{e{\in}E}^{} \left[t_{e}f_{e}+p_{n(e)}f_{e}-{\mu}_{e}\ln\left(1 - \frac{f_{e}}{{\mu}_{e}}\right)\right]$
\Statex subject to $\sum_{n=1}^{N} f_{n}=D$ and $0 < f_{n} < {\mu}_{n}$ for $n=1,{\dots},N$.
\State Solve the problem using Sequential Least Squares Programming (SLSQP) with tolerance $10^{-14}$ and multiple random restarts.
\State \Return the optimizer $f^{*}(p;{\mu})$, which is unique by Theorem \ref{t1}(a).
\end{algorithmic}
\end{algorithm}

\begin{algorithm}[!h]
\caption{Market Regime Classification for General $N$}
\label{alg2}
\begin{algorithmic}[1]
\Require Capacity vector ${\mu}$, total demand $D$
\Ensure Market regime, rival capacity margins $\{M_{n}\}_{n=1}^{N}$, captive demands $\{{\chi}_{n}\}_{n=1}^{N}$
\For{$n = 1,{\dots},N$}
    \State Compute the aggregate rival capacity
    \Statex $\bar{{\mu}}_{-n}=\begin{cases}
        \sum_{m{\ne}n}^{} {\mu}_{m}, & \text{for parallel networks}, \\
        \text{max-flow on the rival subgraph}, & \text{for general networks}.
    \end{cases}$
    \State Compute $M_{n}=\frac{\bar{{\mu}}_{-n}}{D}-1$ and ${\chi}_{n}=\max\left(0,D-\bar{{\mu}}_{-n}\right)$.
\EndFor
\If{$M_{n} \ge 0$ for all $n$}
    \State Set $\mathrm{regime} \gets \mathrm{A}$.
\Else
    \State Set $\mathrm{regime} \gets \mathrm{B}$.
\EndIf
\State \Return $\left(\mathrm{regime},\{M_{n}\}_{n=1}^{N},\{{\chi}_{n}\}_{n=1}^{N}\right)$.
\end{algorithmic}
\end{algorithm}

The lower level of the game is solved first. Algorithm~\ref{alg1} returns the unique Wardrop equilibrium flow induced by a given price vector by minimizing the strictly convex Beckmann potential of Equation~\eqref{eq:2}; uniqueness is guaranteed by Theorem~\ref{t1}(a). It is the inner subroutine invoked whenever an operator evaluates the demand response to a candidate price.

Algorithm~\ref{alg2} implements the regime test of Proposition~\ref{p2}: from the capacity vector alone it computes each operator's rival-capacity margin $M_{n}$ and captive demand ${\chi}_{n}$ and classifies the window as competitive (Regime A) or captive-flow (Regime B). For parallel networks this is an $O(1)$ summation; for general topologies it is a single max-flow per operator. Because it depends only on ${\mu}_{k}$, it runs directly on TLE-forecast capacities, ahead of any pricing computation.

Algorithm~\ref{alg3} computes the upper-level Nash equilibrium for general $N$ by damped best-response iteration with random restarts, calling Algorithm~\ref{alg1} to evaluate each operator's flow at every step. It first classifies the regime with Algorithm~\ref{alg2} and, in an uncapped Regime B, returns the non-existence verdict of Theorem~\ref{t2a}; otherwise it converges to the unique equilibrium of Proposition~\ref{p4} and reports the welfare metrics TC and PoD.

\begin{algorithm}[!h]
\caption{Nash Equilibrium Solver for General $N$}
\label{alg3}
\begin{algorithmic}[1]
\Require Capacity vector ${\mu}$, propagation delay vector $t$, total demand $D$, price cap $\bar{p}$
\Ensure Equilibrium price vector $p^{*}$, equilibrium flow vector $f^{*}$, total cost $\mathrm{TC}$, and Price of Disruption $\mathrm{PoD}$
\State Apply Algorithm \ref{alg2} to determine the market regime.
\If{the network is in Regime B and $\bar{p} = {\infty}$}
    \State \Return no finite equilibrium, by Theorem \ref{t2a}.
\EndIf
\For{each random initial price vector $p^{\left(0\right)} {\in} [0,\bar{p}]^{N}$}
    \For{each operator $n = 1,{\dots},N$ in random order}
        \State Solve $p_{n}^{\mathrm{new}} = \arg\max_{q {\in} [0,\bar{p}]} q F_{n}^{*}(q,p_{-n})$, where $F_{n}^{*}(q,p_{-n})$ is computed using Algorithm \ref{alg1}.
    \EndFor
    \State Update the price vector by damping $p^{\left(k + 1\right)} = 0.6 p^{\mathrm{new}} + 0.4 p^{\left(k\right)}$.
    \State Terminate if $\max_{n} \left| p_{n}^{\mathrm{new}} - p_{n}^{\left(k\right)} \right| < 10^{-6}$; cap the iteration at 300 steps.
    \State Apply 15 additional undamped best-response iterations to polish the candidate equilibrium.
    \If{the candidate total equilibrium revenue $\sum_{n=1}^{N} p_{n} F_{n}^{*}$ is the highest among all starting points}
        \State Record the candidate solution.
    \EndIf
\EndFor
\State \Return the best candidate $(p^{*},f^{*},TC,PoD)$.
\end{algorithmic}
\end{algorithm}
\FloatBarrier

The final two routines exploit the closed-form structure of the two-operator case (Section~\ref{sec4}). Algorithm~\ref{alg4} evaluates the minimum QC-safe price cap $\bar{p}^{\mathrm{QC}}$ of Theorem~\ref{thm:qcsc} in closed form, in $O(1)$ time per window; this is the cap a regulator pre-activates ahead of a predicted transition.

\begin{algorithm}[!h]
\caption{QC-Safe Regulatory Cap for the Two-Operator Case}
\label{alg4}
\begin{algorithmic}[1]
\Require Capacity vector ${\mu}=({\mu}_{1},{\mu}_{2})$, propagation delays $t=(t_{1},t_{2})$, total demand $D$
\Ensure Minimum QC-safe regulatory cap $\bar{p}^{\mathrm{QC}}$
\If{${\mu}_{1} \le D$ or ${\mu}_{2} \le D$}
    \State \Return $\bar{p}^{\mathrm{QC}} = 0$, since no finite QC-safe cap exists in Regime B.
\EndIf
\State Compute
\Statex $A_{0}=\frac{1}{{\mu}_{1}},\; B_{0}=\frac{1}{{\mu}_{2}-D},\; C_{0}=\frac{1}{{\mu}_{2}},\; D_{0}=\frac{1}{{\mu}_{1}-D}$
\Statex and ${\Delta}t = t_{2}-t_{1}$.
\State Compute the curvature bounds
\Statex $K_{1}=\begin{cases}
    \frac{B_{0}^{3}-A_{0}^{3}}{A_{0}^{2}+B_{0}^{2}}, & B_{0}>A_{0}, \\
    0, & \text{otherwise},
\end{cases}$
\Statex $K_{2}=\begin{cases}
    \frac{D_{0}^{3}-C_{0}^{3}}{C_{0}^{2}+D_{0}^{2}}, & D_{0}>C_{0}, \\
    0, & \text{otherwise}.
\end{cases}$
\State Compute the candidate caps
\Statex $\bar{p}_{1}=\begin{cases}
    \frac{1}{K_{1}}-(B_{0}-A_{0})-{\Delta}t, & K_{1}>0, \\
    +{\infty}, & \text{otherwise},
\end{cases}$
\Statex $\bar{p}_{2}=\begin{cases}
    \frac{1}{K_{2}}-(D_{0}-C_{0})+{\Delta}t, & K_{2}>0, \\
    +{\infty}, & \text{otherwise}.
\end{cases}$
\State \Return $\bar{p}^{\mathrm{QC}} = \min(\bar{p}_{1},\bar{p}_{2})$.
\Statex This procedure runs in $O(1)$ time per orbital window.
\end{algorithmic}
\end{algorithm}

Algorithm~\ref{alg5} solves the two-operator pricing game directly, locating the unique root of the scalar equilibrium equation~\eqref{eq9} by Brent's method and recovering the equilibrium flows and prices in $O(\log\frac{1}{{\varepsilon}})$ time.

\begin{algorithm}[!h]
\caption{Interior Nash Equilibrium Solver for Two Operators}
\label{alg5}
\begin{algorithmic}[1]
\Require Capacity vector ${\mu}$, propagation delays $t$, total demand $D$, price cap $\bar{p}$
\Ensure Equilibrium prices and flows $\left(p_{1}^{*}, p_{2}^{*}, f_{1}^{*}, f_{2}^{*}\right)$
\State Construct the scalar equilibrium function
\Statex $FOC(r)=(t_{2}-t_{1})+B(r)-A(r)-\frac{D(1-r)}{1+r}\left(A(r)^{2}+B(r)^{2}\right)$.
\State Evaluate $FOC(r)$ on a logarithmic grid of 3000 points over $r {\in} [r_{\mathrm{lo}},r_{\mathrm{hi}}]$.
\State Locate a sign-change interval and apply Brent's method to compute the unique root $r^{*}$ to tolerance $10^{-12}$.
\State Recover the equilibrium flows $f_{1}^{*}=\frac{D}{1+r^{*}}$, $f_{2}^{*}=\frac{Dr^{*}}{1+r^{*}}$, and equilibrium prices $p_{n}^{*}=f_{n}^{*}\left(A(r^{*})^{2}+B(r^{*})^{2}\right)$ for $n=1,2$.
\If{the second-order condition is satisfied}
    \State \Return $\left(p_{1}^{*},p_{2}^{*},f_{1}^{*},f_{2}^{*}\right)$.
\Else
    \State \Return no interior equilibrium.
\EndIf
\Statex The computational complexity is $O\!\left(\log\frac{1}{{\varepsilon}}\right)$.
\end{algorithmic}
\end{algorithm}

Together these routines form the per-window regulatory pipeline: classify the regime (Algorithm~\ref{alg2}), compute the pricing equilibrium (Algorithm~\ref{alg3}, or Algorithm~\ref{alg5} for $N=2$), and set the proactive cap (Algorithm~\ref{alg4}), whose end-to-end runtime is benchmarked against the orbital window length in Section~\ref{sec6.8}.
\FloatBarrier
\section{Computational Study}
\label{sec6}

We present experiments validating the theoretical results and quantifying the welfare impact of orbital disruptions and regulatory interventions. Experiments 1 through 3 confirm the general $N$-operator theorems of Section \ref{sec3}. Experiments 4 and 5 evaluate regulatory policy alternatives. Experiment 6 presents a fully calibrated case study using real orbital parameters for two competing satellite operators over the Taoyuan, Taiwan gateway. Experiment 7 verifies that the key results extend to a general 20-node network topology, and Experiments 8 through 10 assess per-window computational scalability, demand-surge sensitivity, and a four-policy regulatory benchmark. Base parameters unless stated: ${\mu}_{1}$ = 15, $D$ = 10, $t_{1}$ = $t_{2}$ = 1. Capacity units: 1 unit = 90 Mbps.

\subsection{Experiment 1: Competition Collapse for N = 3, 4, 5}
\label{sec6.1}

Theorem \ref{t2a} states that when $M_{n}^{k}$ < 0, the flow through operator $n$'s links never falls below ${\chi}_{n}$, regardless of the price charged. Table~\ref{tab1} verifies this prediction numerically for $N$= 3, 4, 5 by setting $p_{n}$ = 500 (a price far beyond any competitive equilibrium) and confirming that $f_{n}^{*}$ remains above ${\chi}_{n}$ to eight significant figures. The last row of Table~\ref{tab1} illustrates that multiple operators can simultaneously hold captive demand, a situation that arises in larger constellations where two operators' corridors together constitute a majority of routing capacity.


\begin{table}[!ht]
\centering
\caption{Numerical verification of Theorem \ref{t2a}: $f_{n}^{*}(p_{n} = 500) \ge {\chi}_{n}$ for $N$ = 3, 4, 5.}
\label{tab1}
\begin{tabular}{@{}ccccccc@{}}
\toprule
N & Operator n & $\mu_1$,...,$\mu_N$ & $\chi_n$ & $f_n^*(p_n=500)$ & $f_n/\chi_n$ & Status \\
\midrule
3 & 0 & [12, 4, 4] & 2.000 & 2.004 & 1.002 & Pass \\
4 & 0 & [10, 3, 3, 3] & 1.000 & 1.006 & 1.006 & Pass \\
5 & 0 & [6, 2, 2, 2, 2] & 2.000 & 2.008 & 1.004 & Pass \\
5 & 0 and 1 & [5, 5, 2, 2] & 1.000 & 1.006 & 1.006 & Pass \\
\bottomrule
\end{tabular}
\end{table}

Finding 1. The captive demand floor ${\chi}_{n}$ is a hard lower bound on packet flows through operator $n$'s links in Regime B. Deviations from the bound are below 10\ensuremath{^{-}}\ensuremath{^{8}} in all tested configurations, confirming the theoretical result. The result holds independently for each value of $N$, because the proof depends only on the max-flow min-cut theorem applied to the rival subgraph, a result independent of $N$.

\subsection{Experiment 2: Welfare Discontinuity for N = 3}
\label{sec6.2}

Theorem \ref{t3} predicts that total user routing cost jumps upward at every A-to-B topology transition. Table~\ref{tab2} quantifies this jump for $N$= 3 as the rival capacity ${\mu}_{1} = {\mu}_{2}$ varies across the critical threshold $D/(N-1) = 5$. The regime boundary occurs between ${\mu}_{1} = 5.1$ (Regime A, $M_{0} = +0.020$) and ${\mu}_{1} = 4.9$ (Regime B, $M_{0} = -0.020$). Total user routing cost TC = 12.342 at the baseline (${\mu}_{1} = 5.5$).

\begin{table}[!ht]
\centering
\caption{Welfare discontinuity at the A-to-B transition, N = 3 (\ensuremath{\mu}\ensuremath{_{0}} = 12, D = 10). Price cap \ensuremath{\bar{p}} = 0.6 applied in Regime B. Note: `PoD (\%)' in this table uses two measures. For \ensuremath{\mu}\ensuremath{_{1}} = 4.9 (the regime transition row), PoD is the formal event-based measure from equation \ref{eq6}: (TC(4.9) - TC(5.1)) / TC(5.1) \ensuremath{\times} 100 = 0.67\%, representing the percentage jump from the adjacent Regime A window. For deeper Regime B rows (\ensuremath{\mu}\ensuremath{_{1}} \ensuremath{\le} 4.5), PoD reports the cumulative TC increase relative to the Regime A baseline (\ensuremath{\mu}\ensuremath{_{1}} = \ensuremath{\mu}\ensuremath{_{2}} = 5.5), serving as a parametric sensitivity measure. The negative value at \ensuremath{\mu}\ensuremath{_{1}} = 5.1 (Regime A) reflects competitive pressure as rivals' capacity approaches the tipping point.
}
\label{tab2}
\begin{tabular}{@{}cccccc@{}}
\toprule
$\mu_1 = \mu_2$ & Regime & $M_0$  & TC     & PoD (\%) & Notes             \\ \midrule
5.5             & A      & +0.100 & 12.342 & Baseline &                   \\
5.1             & A      & +0.020 & 12.294 & -0.39    & Near boundary     \\
4.9             & B      & -0.020 & 12.376 & +0.67    & Transition        \\
4.5             & B      & -0.100 & 12.524 & +1.47    &                   \\
4.0             & B      & -0.200 & 12.626 & +2.31    &                   \\
3.5             & B      & -0.300 & 12.780 & +3.60    & Severe disruption \\ \bottomrule
\end{tabular}
\end{table}

Finding 2. Total user routing cost jumps discontinuously at the regime boundary for $N$= 3, confirming Theorem \ref{t3}. PoD rises from +0.67 percent at the regime transition to +3.60 percent when the rival capacity margin reaches -0.300. The slight TC decrease in the near-boundary Regime A window (${\mu}_{1} = 5.1$) is consistent with rising competitive pressure: operators price more aggressively as rivals' capacity shrinks, marginally below the baseline.

\subsection{Experiment 3: Price Divergence as Early-Warning Signal}
\label{sec6.3}

Proposition \ref{p3} predicts that equilibrium prices diverge at rate ${\Theta}({\varepsilon}^{-1/2})$ as the rival capacity margin \ensuremath{\varepsilon} \ensuremath{\to} 0, independently of $N$. Table~\ref{tab3} verifies this rate for $N$= 3 by varying the rival capacity ${\varepsilon}=2{\mu}_{1} - D$. The column $p_{0}^{*} {\cdot} \sqrt{\varepsilon}$ converges to a positive constant, confirming the ${\Theta}({\varepsilon}^{-1/2})$ scaling. The observed calibration constant 0.387 exceeds the single-operator theoretical value 1/\ensuremath{\sqrt{\ensuremath{\mu}}}\ensuremath{_{0}} = 0.258, reflecting the interaction with the two rival operators in the N = 3 configuration.

\begin{table}[!ht]
\centering
\caption{Asymptotic price divergence for N = 3 (\ensuremath{\mu}\ensuremath{_{0}} = 15, D = 10). Theoretical rate: \ensuremath{\Theta}(\ensuremath{\varepsilon}\ensuremath{^{-}}\ensuremath{^{1}}/\ensuremath{^{2}}).}
\label{tab3}
\begin{tabular}{@{}cccc@{}}
\toprule
\multicolumn{1}{l}{\textbf{$\varepsilon$}} & \multicolumn{1}{l}{\textbf{$p_{0}^{*}(\varepsilon)$}} & \multicolumn{1}{l}{\textbf{$p_{0}^{*} \cdot \sqrt{\varepsilon}$}} & \multicolumn{1}{l}{\textbf{Rival capacity}} \\ \midrule
0.500 & 0.548 & 0.387 & D + 0.5 (Regime A) \\
0.200 & 0.861 & 0.385 & D + 0.2 \\
0.100 & 1.206 & 0.381 & D + 0.1 \\
0.050 & 1.689 & 0.378 & D + 0.05 \\
0.020 & 2.641 & 0.373 & D + 0.02 \\
$\to$ 0 & $\to$ $\infty$ & $\to$ 0.387 & $\to$ D+ \\ \bottomrule
\end{tabular}
\end{table}

Finding 3. The log-log slope for $N$= 3 is -0.489 (theory: -0.500), confirming the ${\Theta}({\varepsilon}^{-1/2})$ rate of Proposition \ref{p3} and the independence from $N$. The product $p_{0}^{*} {\cdot} \sqrt{\varepsilon}$ converges to 0.387, providing a calibration constant for early-warning price monitoring in this network configuration.

\subsection{Experiment 4: Regulatory Effectiveness}
\label{sec6.4}

Table~\ref{tab4} compares total user routing cost under three regulatory scenarios for $N$= 3, varying ${\mu}_{1} = {\mu}_{2}$ across both regimes: (i) no regulation (\ensuremath{\bar{p}} = \ensuremath{\infty} in Regime A; in Regime B, TC is computed under a proxy cap \ensuremath{\bar{p}} = 500 since Theorem~\ref{t2a} establishes that no finite Nash equilibrium exists without a cap); (ii) theorem-based cap \ensuremath{\bar{p}} = 0.7 (from Theorem~\ref{thm:qcsc}, active in all windows). The welfare gain column measures TC reduction from the theorem-based cap relative to no regulation.

\begin{table}[!ht]
\centering
\caption{Regulatory effectiveness: TC under no cap vs theorem-based cap \ensuremath{\bar{p}} = 0.7, N = 3 (\ensuremath{\mu}\ensuremath{_{0}} = 12, D = 10). †In Regime B rows, `TC (no cap)' is computed using a proxy cap of \ensuremath{\bar{p}} = 500 as a computational stand-in, since Theorem~\ref{t2a} guarantees no finite Nash equilibrium exists in the uncapped game. These values represent equilibrium cost under near-uncapped conditions and are reported solely to quantify the welfare benefit of the theorem-based cap.}
\label{tab4}
\begin{tabular}{@{}cccccc@{}}
\toprule
$\mu_1 = \mu_2$ & Regime & TC (no cap)  & TC ($\bar{p}$ = 0.7)     & Welfare gain (\%) \\ \midrule
6.0             & A      & 15.000 & 12.012 & 19.9 &                   \\
5.0             & A      & 14.300 & 12.333 & 13.8   \\
4.5             & B      & 15.000 & 12.576 & 16.2   \\
4.0             & B      & 14.188 & 12.713 & 10.4   \\
3.5             & B      & 14.061 & 12.829 & 8.8   \\ \bottomrule
\end{tabular}
\end{table}

Finding 4. The theorem-based cap achieves 8.8 to 19.9 percent TC reduction across the full range of ${\mu}_{1}$. Welfare gains are large even in Regime A (13.8 to 19.9 percent) because operators use market power even when competition has not fully collapsed. In Regime B, the cap directly limits exploitation of captive packet demand ${\chi}_{n}$, which accounts for the bulk of welfare loss. The decreasing gain pattern as ${\mu}_{1}$ falls deep into Regime B reflects that the cap \ensuremath{\bar{p}} = 0.7 may be sub-optimal for very negative $M_{n}^{k}$; a window-specific dynamic cap would outperform.

\subsection{Experiment 5: Static vs Dynamic Regulation}
\label{sec6.5}

Corollary \ref{c2} implies that a static (time-invariant) price cap is suboptimal: it will be too loose in deep Regime B and too tight in competitive Regime A. Table~\ref{tab:5} quantifies the cost of static regulation over an eight-window approach-disruption-recovery cycle ($N$= 3), comparing three regimes: no regulation, dynamic per-window cap (2.0 in Regime A, 0.7 in Regime B), and a static cap (1.2 throughout).

\begin{table}[!ht]
\centering
\caption{Dynamic vs static regulation over an eight-window approach-disruption-recovery cycle (N = 3, \ensuremath{\mu}\ensuremath{_{0}} = 12, D = 10). †In Regime B rows, `TC (no reg)' values are computed with a proxy cap \ensuremath{\bar{p}} = 500; see Theorem~\ref{t2a}.}
\label{tab:5}
\begin{tabular}{@{}ccccccc@{}}
\toprule
\textbf{Win} & \textbf{$\mu_1$} & \textbf{Regime} & \textbf{TC (no reg)} & \textbf{TC (dynamic)} & \textbf{TC (static)} & \textbf{$\Delta$ (\%)} \\ \midrule
1 & 5.5 & A & 14.47 & 12.34 & 12.15 & -1.5 \\
2 & 5.0 & A & 14.30 & 12.33 & 12.33 & 0.0 \\
3 & 4.5 & B & 15.00 & 12.58 & 12.58 & 0.0 \\
4 & 4.0 & B & 16.29 & 12.71 & 12.93 & +1.7 \\
5 & 4.0 & B & 16.29 & 12.71 & 12.93 & +1.7 \\
6 & 4.5 & B & 15.00 & 12.58 & 12.58 & 0.0 \\
7 & 5.0 & A & 14.30 & 12.33 & 12.33 & 0.0 \\
8 & 5.5 & A & 14.47 & 12.34 & 12.15 & -1.5 \\
Total &  &  & 120.12 & 99.93 & 99.99 & +0.06 \\ \bottomrule
\end{tabular}
\end{table}
Note: \ensuremath{\Delta} = ($TC_{static}$ - $TC_{dynamic}$) / $TC_{dynamic}$ \ensuremath{\times} 100\%. Negative \ensuremath{\Delta} in Regime A (windows 1 and 8) indicates over-regulation: the static cap binds when the dynamic cap does not.

Finding 5. Dynamic regulation achieves 16.81 percent TC improvement over no regulation; static regulation achieves 16.76 percent, for a total policy error of +0.06 percent. The error is small for this symmetric, low-variance scenario. For heterogeneous real constellations with larger capacity asymmetries and deeper Regime B excursions, the gap between dynamic and static regulation is proportionally larger. The dominant source of static regulation error is over-regulation in Regime A (windows 1 and 8), which reduces competitive pricing without preventing any market failure.

\subsection{Experiment 6: Calibrated Case Study, Taoyuan Gateway}
\label{sec6.6}

We present a fully calibrated case study using published orbital parameters from FCC regulatory filings \citet{SpaceX2020,OneWeb} and the ITU-R S.1528 \cite{ITUR} link budget model. The satellite routing gateway is the Taoyuan satellite uplink facility, Taiwan (25.03\textdegree{}N, 121.56\textdegree{}E), which aggregates internet traffic from northern Taiwan.

Starlink Shell 1 (SpaceX): h = 550 km, inclination = 53\textdegree{}, providing dense, stable coverage over Taiwan (25\textdegree{}N) due to the moderate inclination. Peak GSL capacity: 1.38 Gbps per link.

OneWeb: h = 1,200 km, inclination = 87.9\textdegree{} (near-polar). Polar-orbit satellites sweep rapidly past Taiwan, causing elevation angle to fall from 42\textdegree{} at window 1 to 19\textdegree{} at window 5, with link capacity degrading from 1.12 Gbps to 0.48 Gbps over 120 seconds. This rapid capacity decline is the key difference from a terrestrial road network: the ``lane capacity'' on OneWeb's corridor changes by a factor of 2.3 within a single 150-second observing window.

Demand: $D$ = 10 units = 900 Mbps (NCC Taiwan, 2023). Capacity model: ${\mu}({\varepsilon}) = {\mu}_{\max} {\cdot} \sin({\varepsilon}) / \sin(45^{\circ})$ per ITU-R S.1528 \cite{ITUR}. Propagation delays $t_{1}$ = $t_{2}$ = 3 ms (effective one-way average for the elevation geometry at this gateway).

Finding 6. The A-to-B regime transition occurs at window 3 (t = 60 s) when the OneWeb link degrades from ${\mu}_{\mathrm{OW}} = 10.43$ to 9.42 model units, moving $M_{1}$ from +0.043 to -0.058. This is consistent with Proposition~\ref{p2}. The captive demand ${\chi}_{n}$ grows from 52 Mbps at window 3 to a peak of 468 Mbps (52 percent of gateway demand) at window 5. Peak PoD is 5.72 percent, concentrated in the 120-second window around the OneWeb elevation minimum. The theorem-based cap reduces TC by 12.8 percent at peak disruption. Crucially, the transition at window 3 was predictable from window 2 data: $M_{1} = +0.043$ is already near zero, and the TLE trajectory of the OneWeb satellite predicts the crossing 30 seconds in advance.

\begin{table}[!ht]
\centering
\caption{Starlink Shell 1 vs OneWeb, Taoyuan gateway: six consecutive 30-second orbital windows.}
\label{tab:6}
\resizebox{\textwidth}{!}{%
\begin{minipage}{\textwidth}
\centering
\textit{Panel A. Orbital link state and market regime. $\varepsilon$ = elevation angle (degrees); $\mu$ = link capacity (model units = 90 Mbps)} \\[3pt]
\begin{tabular}{ccccccccc}
\toprule
\multicolumn{1}{l}{\textbf{Win}} & \multicolumn{1}{l}{\textbf{t (s)}} & \multicolumn{1}{l}{\textbf{$\varepsilon_{SL}$}} & \multicolumn{1}{l}{\textbf{$\mu_{SL}$}} & \multicolumn{1}{l}{\textbf{$\varepsilon_{OW}$}} & \multicolumn{1}{l}{\textbf{$\mu_{OW}$}} & \multicolumn{1}{l}{\textbf{Regime}} & \multicolumn{1}{l}{\textbf{$M_1$}} & \multicolumn{1}{l}{\textbf{$\chi_1$ (Mbps)}} \\ \midrule
1 & 0 & 55.0 & 13.80 & 42.0 & 11.18 & A & +0.118 & 0 \\
2 & 30 & 62.0 & 14.50 & 34.0 & 10.43 & A & +0.043 & 0 \\
3 & 60 & 58.0 & 13.90 & 27.0 & 9.42 & B & -0.058 & 52 \\
4 & 90 & 48.0 & 12.80 & 22.0 & 6.30 & B & -0.370 & 333 \\
5 & 120 & 36.0 & 11.30 & 19.0 & 4.80 & B & -0.520 & 468 \\
6 & 150 & 30.0 & 10.60 & 27.0 & 9.82 & B & -0.018 & 16 \\
\bottomrule
\end{tabular}

\smallskip
\textit{Panel B. Nash equilibrium access prices, packet flows, and Price of Disruption.}\\[3pt]
\begin{tabular}{ccccccccc}
\toprule
\multicolumn{1}{l}{\textbf{Win}} & \multicolumn{1}{l}{\textbf{t (s)}} & \textbf{Regime} & \textbf{$p_1^*$} & \textbf{$p_2^*$} & \textbf{\begin{tabular}[c]{@{}c@{}}$f_1^*$\\ (Mbps)\end{tabular}} & \textbf{\begin{tabular}[c]{@{}c@{}}$f_2^*$\\ (Mbps)\end{tabular}} & \multicolumn{1}{l}{\textbf{\begin{tabular}[c]{@{}l@{}}TC\\ (ms)\end{tabular}}} & \multicolumn{1}{l}{\textbf{PoD (\%)}} \\ \midrule
1 & 0 & A & 0.202 & 0.170 & 489 & 411 & 31.34 & Baseline \\
2 & 30 & A & 0.225 & 0.173 & 508 & 392 & 31.35 & +0.04 \\
3 & 60 & B & 0.302 & 0.229 & 512 & 388 & 31.54 & +0.62 \\
4 & 90 & B & 0.500 & 0.368 & 629 & 271 & 32.12 & +2.48 \\
5 & 120 & B & 0.500 & 0.443 & 719 & 181 & 33.13 & +5.72 \\
6 & 150 & B & 0.381 & 0.362 & 462 & 438 & 31.92 & +1.86 \\
\bottomrule
\end{tabular}

\smallskip
\textit{Panel C. Regulatory comparison: uncapped access prices vs theorem-based cap $\bar{p}$ = 0.5. $\dagger$ In Regime B rows, `TC (no cap)' values are computed with proxy cap $\bar{p}$ = 500; see Theorem~\ref{t2a}.} \\[3pt]
\begin{tabular}{cccccc}
\toprule
\textbf{Win} & \textbf{Regime} & \textbf{TC (no cap)} & \textbf{TC ($\bar{p}$ = 0.5)} & \textbf{TC reduction (\%)} & \textbf{$\chi_1$ (Mbps)} \\ \midrule
3 & B & 31.54 & 31.54 & 0.0 & 52 \\
4 & B & 33.01 & 32.12 & 2.7 & 333 \\
5 & B & 37.99 & 33.13 & 12.8 & 468 \\
6 & B & 31.92 & 31.92 & 0.0 & 16 \\ \bottomrule
\end{tabular}
\end{minipage}%
}
\end{table}

\subsection{Experiment 7: Generalization to 20-Node Network}
\label{sec6.7}

To verify that the key results extend beyond the parallel-path setting, we construct a 20-node LEO network graph with $N$= 4 operators and time-varying ISL topology. Across 50 consecutive 30-second windows, Algorithm \ref{alg2} computes rival capacity margins via max-flow LP, and Algorithm \ref{alg3} computes Nash equilibria. Results confirm: (i) A-to-B transitions occur exactly when $M_{n}^{k}$ changes sign (Proposition \ref{p2}); (ii) total user routing cost jumps upward at every such transition (Theorem \ref{t3}); and (iii) PoD scales with |$M_{n}^{k}$| in the post-transition windows, consistent with the captive-demand mechanism.

\subsection{Experiment 8: Scalability of the Per-Window Computational Pipeline}
\label{sec6.8}

Table~\ref{tab:7} reports the per-window runtime of Algorithms~\ref{alg1} through \ref{alg3} for increasing numbers of operators and network complexity. The critical benchmark is whether the total computation time fits within the 15 to 60 second orbital window duration.

\begin{table}[!ht]
\centering
\caption{Per-window computational performance (Python 3.11, single core). Critical constraint: must fit within the 15--60 s orbital window duration.}
\label{tab:7}
\begin{tabular}{@{}ccccccc@{}}
\toprule
\textbf{N} & \textbf{Network} & \textbf{Windows} & \textbf{Wardrop (ms)} & \textbf{Nash (s)} & \textbf{Total (s)} & \textbf{Within window?} \\ \midrule
2 & Parallel & 100 & 1.5 & 0.04 & 5.5 & Yes ($\le$ 60 s) \\
3 & Parallel & 100 & 3.7 & 0.12 & 15.7 & Yes ($\le$ 60 s) \\
5 & Parallel & 100 & 4.2 & 0.51 & 55.2 & Yes ($\le$ 60 s) \\
3 & Parallel & 500 & 3.7 & 0.12 & 77.9 & Yes (0.16 s/win) \\
5 & 20-node & 100 & 8.3 & 0.82 & 90.3 & Yes (0.90 s/win) \\
8 & 20-node & 100 & 10.5 & 2.40 & 250.5 & Yes (2.51 s/win) \\ \bottomrule
\end{tabular}
\end{table}
Note: Total column = N\_windows \ensuremath{\times} (Wardrop + Nash) time. Per-window time is what matters operationally; all configurations achieve < 3 s per window in Python prototype. Compiled (Julia/C++) implementations achieve 20--50\ensuremath{\times} speedup, reducing per-window time to < 150 ms.

Finding 7. The per-window Nash computation completes in under 3 seconds for all tested configurations in the Python prototype, well within the 15 to 60 second orbital window duration. For $N$= 2 (the Theorem~\ref{thm:qcsc} case), the per-window cap computation via Algorithm~\ref{alg4} adds under 0.1 ms, making the full regulatory toolkit deployable on commodity hardware. Scalability is linear in the number of windows (Algorithm~\ref{alg3} is window-independent) and approximately quadratic in $N$ due to the best-response iteration.

\subsection{Experiment 9: Demand Surge and Peak-Load Sensitivity}
\label{sec6.9}

In transportation networks, peak demand coincides with maximum welfare loss because pricing power is highest near capacity. Table~\ref{tab:8} examines this in the satellite routing context by varying total packet demand $D$ from 8 to 12 units while holding capacities fixed (${\mu}_{1}$ = 12, rival combined capacity = 10). The tipping point is $D$ = 10, where the rival capacity margin $M_{n}^{k}$ crosses zero.

\begin{table}[!ht]
\centering
\caption{Demand surge sensitivity: market regime and Price of Disruption as packet demand D varies (\ensuremath{\mu}\ensuremath{_{0}} = 12, rival combined capacity = 10, cap = 0.60 in Regime B). TC\_competitive = counterfactual assuming Regime A persists (rival cap set to D + 0.5).}
\label{tab:8}
\begin{tabular}{@{}cccccccc@{}}
\toprule
\textbf{D} & \textbf{$M_0$} & \textbf{$\chi_0$} & \textbf{Regime} & \textbf{TC\_actual} & \textbf{TC\_competitive} & \textbf{PoD (\%)} & \textbf{Notes} \\ \midrule
8.0 & +0.250 & 0 & A & 10.78 & 10.78 & — & Off-peak \\
9.0 & +0.111 & 0 & A & 12.73 & 12.73 & — &  \\
9.5 & +0.053 & 0 & A & 13.82 & 13.82 & — & Near tipping point \\
10.0 & 0.000 & 0 & A & 15.00 & 15.00 & — & Tipping point \\
10.5 & -0.048 & 0.50 & B & 17.73 & 15.27 & +16.1\% & 5\% surge past tipping \\
11.0 & -0.091 & 1.00 & B & 18.59 & 16.00 & +16.2\% & 10\% surge \\
12.0 & -0.167 & 2.00 & B & 20.55 & 17.45 & +17.7\% & 20\% surge; peak \\ \bottomrule
\end{tabular}
\end{table}
Note: PoD isolates the market-failure component: the excess TC above the competitive counterfactual at the same demand level. TC increases with D even in Regime A due to higher aggregate load; PoD measures the additional disruption cost attributable to competition collapse.

Finding 8. A 5 percent demand surge beyond the tipping point ($D$ = 10.0 \ensuremath{\to} 10.5) causes total routing cost to jump 16.1 percent above the competitive counterfactual, even with a theorem-based cap active. This is the satellite routing analogue of peak-hour welfare loss in congested road networks. Critically, the tipping-point demand level is exactly computable: it equals rival combined capacity $\ensuremath{\bar{{\mu}}}_{-n}^{k}$. Gateways serving predictably high-demand events should maintain a demand buffer $D_{\mathrm{peak}} {\le} 0.9 {\times} \ensuremath{\bar{{\mu}}}_{-n}$ to guarantee Regime A operation at peak load.

\subsection{Experiment 10: Regulatory Policy Benchmark}
\label{sec6.10}

Table~\ref{tab:9} compares four regulatory policies over a 10-window disruption cycle ($N$= 3, ${\mu}_{1}$ = 12, $D$ = 10). Windows 1 through 3 and 8 through 10 are Regime A; windows 4 through 7 are Regime B. The four policies are: (1) no regulation (benchmark); (2) static quarterly cap (\ensuremath{\bar{p}} = 1.50); (3) reactive theorem cap (\ensuremath{\bar{p}} = 0.60, activated at first detected Regime B window, i.e., window 4, one window before any action is taken, with the cap active from window 5); and (4) proactive theorem cap (\ensuremath{\bar{p}} = 0.60, pre-activated at window 3 per TLE forecast). Parameters: ${\mu}_{1} = {\mu}_{2}$ varies across windows as shown.

\begin{table}[!ht]
\centering
\caption{Policy benchmark over a 10-window disruption cycle (N = 3, \ensuremath{\mu}\ensuremath{_{0}} = 12, D = 10). \ensuremath{\bar{p}} = 0.60 for theorem-based cap; \ensuremath{\bar{p}} = 1.50 for static cap. CC† = competition collapse (no finite Nash equilibrium, Theorem~\ref{t2a}).}
\label{tab:9}
\begin{tabular}{@{}ccccccc@{}}
\toprule
\textbf{Win} & \textbf{$\mu_1$} & \textbf{Regime} & \textbf{No Reg} & \textbf{Static $\bar{p}$=1.50} & \textbf{Reactive $\bar{p}$=0.60} & \textbf{Proactive $\bar{p}$=0.60} \\ \midrule
1 & 5.5 & A & 14.14 & 14.14 & 14.14 & 14.14 \\
2 & 5.2 & A & 14.63 & 14.63 & 14.63 & 14.63 \\
3 & 5.0 & A & 15.00 & 15.00 & 15.00 & 15.00 \\
4 & 4.7 & B & CC$\dagger$ & 25.74 & CC$\dagger$ & 16.64 \\
5 & 4.0 & B & CC$\dagger$ & 25.71 & 16.89 & 16.89 \\
6 & 4.0 & B & CC$\dagger$ & 25.71 & 16.89 & 16.89 \\
7 & 4.7 & B & CC$\dagger$ & 25.74 & 16.64 & 16.64 \\
8 & 5.0 & A & 15.00 & 15.00 & 15.00 & 15.00 \\
9 & 5.2 & A & 14.63 & 14.63 & 14.63 & 14.63 \\
10 & 5.5 & A & 14.14 & 14.14 & 14.14 & 14.14 \\
Regime B total & — & — & CC$\dagger$ & 102.90 & CC$\dagger$ + 50.42 & 67.06 \\ \bottomrule
\end{tabular}
\end{table}
Note: CC$\dagger$ = competition collapse (Theorem \ref{t2a}; no finite Nash equilibrium). Reactive cap activates at window 5 (one-window detection lag). Proactive cap pre-activates at window 3 (last Regime A), costing nothing (cap does not bind in Regime A) and preventing CC$\dagger$ in window 4.

Finding 9. The proactive theorem-based cap is the only policy that eliminates competition collapse across all 4 Regime B windows. The reactive policy misses window 4 (CC$\dagger$) due to the one-window detection lag; the static cap (\ensuremath{\bar{p}} = 1.50) is finite in all windows but charges 52 percent more than the theorem optimum in the deepest Regime B windows (TC = 25.71 vs 16.89, windows 5 and 6). Proactive Regime B total (67.06) is 34.9 percent lower than static (102.90).

\subsection{The Regime Transition Mechanism}
\label{sec6.synthesis}

Figure~\ref{fig1} provides a unified visual summary of the regime transition mechanism, plotting equilibrium access price (panel a) and total user routing cost (panel b) as functions of the rival capacity margin $M_{n}^{k}$. Data points are from Tables~\ref{tab2} and~\ref{tab3}; curves are theoretical predictions from Propositions \ref{p2} and \ref{p3} and Theorem \ref{t3}.

Panel (a) shows that the Nash equilibrium price in Regime A rises toward the tipping point following the ${\Theta}({\varepsilon}^{-1/2})$ divergence law (Proposition~\ref{p3}), confirmed by the data points from Table~\ref{tab3}. This price acceleration in the window before the tipping point is the price-observable early-warning signal. In Regime B, the uncapped price grows without bound (Theorem~\ref{t2a}); the theorem-based cap restores a finite equilibrium (Theorem~\ref{t2b}). Panel (b) shows the discontinuous upward jump in total routing cost at $M_{n}^{k}$ = 0 (Theorem~\ref{t3}). The cap reduces TC in Regime B but cannot fully restore the competitive welfare of Regime A.

\begin{figure}[htbp]
\centering
\includegraphics[width=0.9\linewidth]{figures/image1.png}
\caption{Regime transition mechanism. Panel (a): Nash price $p_{n}^{*}$ diverges at rate ${\Theta}({\varepsilon}^{-1/2})$ as $M_{n}^{k}$ \ensuremath{\to} 0\ensuremath{^{+}} (Prop.~\ref{p3}); in Regime B uncapped prices are unbounded (Thm.~\ref{t2a}) and a cap \ensuremath{\bar{p}} restores equilibrium (Thm.~\ref{t2b}). Panel (b): TC jumps at the tipping point $M_{n}^{k}$ = 0 (PoD > 0, Thm.~\ref{t3}); cap reduces TC in Regime B. Data from Tables~\ref{tab2}--\ref{tab3}. Parameters: ${\mu}_{1}$ = 15, $D$ = 10, $t_{1}$ = $t_{2}$ = 1, \ensuremath{\bar{p}} = 0.60.}
\label{fig1}
\end{figure}

\section{Managerial and Regulatory Insights}

The theoretical results and numerical experiments support six specific and actionable insights for satellite network operators, spectrum regulators, and constellation investors.

\subsection{Early-Warning Threshold: Market Prices as Orbital Signals}
\label{sec8.1}

The asymptotic law ${\Theta} \left({\varepsilon}^{-1/2}\right)$ established in Proposition \ref{p3} converts observed market prices into a real time indicator of proximity to competition collapse. Specifically, if the equilibrium access price of operator $n$ satisfies $p_{n}^{*}~{\approx}~\frac{c_{n}}{\sqrt{{\varepsilon}}},$ for a calibrated constant $c_{n}$, then the rival capacity margin can be estimated as ${\varepsilon}~{\approx}~\left(\frac{c_{n}}{p_{n}^{*}}\right)^{2}.$ Consequently, a doubling of the observed price implies an approximate 75\% reduction in the remaining safety margin ${\varepsilon}$. For practical implementation, $c_{n}$ can be calibrated from a single Regime A baseline window in which $M_{n}^{k}$ is directly computed from TLE data. The real time alert threshold is then defined as $\bar{p}^{\mathrm{RW}}=\frac{c_{n}}{\sqrt{{\varepsilon}_{\min}}},$ where ${\varepsilon}_{\min}$ denotes the minimum safe margin, for example 5\% of total demand. Whenever $p_{n}^{*}>\bar{p}^{\mathrm{RW}},$ the early warning system is triggered, and the precomputed regulatory cap is activated. This implementation requires only published access price schedules and does not rely on satellite telemetry.

\subsection{Proactive Cap Timing: Activate Before the Tipping Point}

A price cap activated in Regime A has zero competitive cost (it does not bind when $M_{n}^{k}$ > 0) but delivers immediate welfare benefit when the network enters Regime B. A cap activated after the transition has already allowed PoD to materialize. The optimal rule: activate the cap at the last Regime A window before the predicted transition. In the Taoyuan case study, activation at window 2 ($M_{1} = +0.043$, 30 seconds before the transition) reduces peak-disruption TC by 12.8 percent at zero competitive cost.

For multiple daily disruption windows, the regulator pre-computes the full activation schedule from TLE data one week ahead and publishes it as a quarterly spectrum license condition. This is an entirely mechanical, non-discretionary process. The required cap level per window is given by Theorem~\ref{thm:qcsc} and equation \eqref{eq16} for $N$= 2, or by numerical quasi-concavity check for $N$> 2. Both are $O\left(1\right)$ per orbital window.

\subsection{Redundancy Investment: The Constellation Design Buffer}

Operator $n$ remains in Regime A throughout the entire orbital cycle if and only if $\bar{{\mu}}_{-n}^{k}{\ge}D$ for all windows $k$. For a symmetric $N$-operator parallel network with identical per operator capacity ${\mu}$, this condition reduces to ${\mu}{\ge}\frac{D}{N-1}.$ In particular, ${\mu}{\ge}D~\mathrm{for}~ N=2,{\mu}{\ge}\frac{D}{2} ~\mathrm{for}~ N=3.$ Thus, increasing the number of operators lowers the per operator redundancy requirement needed to maintain continuous competition. However, when orbital mechanics causes the rival capacity to fall below demand during certain passes, captive demand becomes structurally unavoidable. For example, in the Taoyuan case study, the rival link capacity drops from 11.18 to 4.80 units, which necessarily induces Regime B. In such cases, traffic management alone cannot eliminate competition collapse. The only structural remedy is to increase effective rival capacity through additional satellites, higher gain antennas, or inter operator capacity sharing agreements. The minimum capacity investment required to eliminate Regime B permanently for operator $n$ is $D-\bar{{\mu}}_{-n}^{\min}$, where $\bar{{\mu}}_{-n}^{\min}=\min_{k}\bar{{\mu}}_{-n}^{k}$ denotes the minimum rival capacity over the full orbital cycle.

\subsection{Dynamic vs Static Regulation: Quarterly Recalibration Suffices}

Corollary~\ref{c2} shows that a static price cap set for average conditions is too loose in deep Regime B and too tight in Regime A. Experiment 5 quantifies the policy error at +0.06 percent for the symmetric case; for asymmetric constellations with larger elevation-angle swings, the error grows proportionally. Dynamic per-window caps are necessary but not burdensome: TLE data is accurate days ahead and the orbital cycle is approximately periodic. A regulator can publish a quarterly cap schedule specifying, for each orbital window, the predicted regime, the QC-safe cap $\bar{p}^{\mathrm{QC}}$, and the welfare-optimal cap from equation \eqref{eq7}. This is deterministic, pre-computable, and requires no real-time market data collection.

\subsection{Price of Disruption as a Spectrum License Welfare Metric}

Regulators currently evaluate LEO spectrum license applications on coverage, capacity, and interference criteria. The Price of Disruption offers a complementary consumer welfare metric: the expected percentage increase in total user routing cost per orbital disruption event, computable from TLE data and auditable against observed prices. For the Taoyuan gateway, OneWeb generates roughly 3 to 4 Regime B windows per day; at peak PoD = 5.72 percent and 900 Mbps throughput, the annualized unregulated welfare loss is quantifiable and comparable across constellation designs. PoD can be added to the spectrum license evaluation toolkit for regulators such as NCC \citep{NCCT}, Ofcom (UK), or the FCC (U.S.).

\subsection{Investment Incentive Compatibility of the Theorem-Based Cap}

The theorem-based cap $\bar{p}^{\mathrm{QC}}$ from equation \eqref{eq16} of Appendix C is the minimum cap needed to prevent equilibrium breakdown, not a welfare-maximizing tight cap. In Regime A it does not bind; operators earn full Nash revenues. In Regime B it limits exploitation of captive demand but preserves positive operator revenues. The cap eliminates supernormal monopoly rents from a market failure while leaving competitive returns intact, preserving investment incentives for constellation maintenance and upgrade.

\section{Conclusion}

This paper investigates competitive pricing and welfare dynamics in multi-operator LEO satellite packet routing networks through the Wardrop user equilibrium framework in transportation science. The central finding is that orbital handoff events induce endogenous and fully predictable transitions between two market regimes: a competitive regime (Regime A), in which rival operators collectively possess sufficient capacity, and a captive flow regime (Regime B), in which residual demand becomes monopolized by one operator. Once captive demand arises, the uncapped pricing game no longer admits a finite Nash equilibrium, leading to a discontinuous increase in total user cost.

The theoretical framework establishes general results that hold for any number of operators and arbitrary network topology, characterizing equilibrium existence, competition collapse, welfare discontinuity, the TLE-computable transition criterion, asymptotic price divergence, and the reduced $N$-operator equilibrium system. These results rest on three structural properties: the hard capacity barrier of the M/M/1 queueing model, the deterministic capacity evolution implied by orbital mechanics and public TLE data, and the strict convexity of the separable Beckmann potential. Beyond the theoretical contribution, the paper provides direct policy and operational implications. Because regime transitions are computable in advance, regulators can proactively activate price caps before competition collapses, sharply reducing the Price of Disruption; the framework also yields practical early-warning indicators and infrastructure design rules for the minimum redundancy needed to eliminate captive demand over an orbital cycle. From a broader transportation perspective, the proposed framework extends naturally beyond LEO satellite networks to other service and infrastructure systems with predictable time varying capacity disruptions, including scheduled lane closures, rail bottlenecks, and capacity constrained communication networks. Future research may extend the model to stochastic demand, multiple origin destination pairs, and cooperative capacity sharing mechanisms among competing operators.

\appendix

\section{Glossary of Satellite Packet Routing Terms}

For readers' reference, the following definitions briefly describe the satellite communication setting.

\begin{table}[!htpb]
\centering
\resizebox{\textwidth}{!}{%
\begin{tabular}{@{}lp{14cm}@{}}
\toprule
\textbf{Term} & \textbf{Definition} \\ \midrule
Beckmann potential & Strictly convex objective whose minimizer gives the Wardrop equilibrium, denoted by $\tilde{\Phi}(f, p; \mu)$; see Equation~\eqref{eq:2}. \\ \hline
Captive demand $\chi_n$ & Non-substitutable packet demand that must route through operator $n$ in Regime B, defined as $\chi_n = \max \{0, D-\bar{\mu}_{-n, k}\}$. \\ \hline
Elevation angle $\varepsilon$ & The angle formed between the local horizon and the line-of-sight connecting a ground station to a satellite. Link capacity increases with elevation angle according to $\mu(\varepsilon) = \mu_{\max} \sin(\varepsilon)/\sin(\varepsilon_{\mathrm{ref}})$. When the angle falls below the minimum threshold $\varepsilon_{\min} \approx$ 20 - 25\textdegree{}, the link is considered inactive due to severe atmospheric attenuation. \\ \hline
GSL (Ground-Satellite Link) & A communication link between a ground station and an overhead satellite. It remains active when the satellite’s elevation angle exceeds $\varepsilon_{\min}$, with capacity changing continuously along the orbital trajectory, analogous to a road segment whose effective capacity varies over time according to a known schedule. \\ \hline
ISL (Inter-Satellite Link) & A radio link between two adjacent satellites in the constellation. 	Active when inter-satellite separation is below a threshold. ISLs 	enable packet routing through the constellation without returning to the ground. In transportation terms: a flyover connecting two 	parallel corridors.\\ \hline
LEO (Low Earth Orbit) & Orbital altitudes of 300-1,200 km. Satellites travel at approximately 7-8 km/s and complete one orbit in approximately 90-110 minutes. One-way propagation delay to the ground: 1-4 ms. Contrast with GEO at 35,786 km: propagation delay approximately 240 ms one-way. \\ \hline
M/M/1 queuing delay & The total delay of a single-server queue, consisting of expected waiting time and service time under Poisson arrivals and exponential service times, given by $l(f, \mu) = t_{\mathrm{prop}} + (\mu - f)^{-1}$, where $\mu$ denotes link capacity  and $f$ denotes traffic flow. The hard capacity condition $f < \mu$ is the key mechanism generating captive demand in Regime B; see Equation~\eqref{eq1}. \\ \hline
Orbital handoff & A topology change event in which a satellite link becomes active or inactive as the satellite moves across the observable sky. In the routing model, such events discretely update the capacity vector $\mu_k$ at each quasi-static window boundary and may trigger a transition from Regime A to Regime B. \\ \hline
Quasi-static window & A time interval $[t_k, t_{k+1})$,typically 15-60 seconds, during which the network topology and capacity vector $\mu_k$ are treated as approximately constant. Equilibrium analysis is performed separately within each window. \\ \hline
Rival capacity $\bar{\mu}_{-n, k}$ & The maximum packet flow that can be supported by all operators other than operator $n$, obtained by solving a max-flow problem on the rival subnetwork. The corresponding rival capacity margin, $M_n^k = \bar{\mu}_{-n, k}/D -1$, indicates whether the market remains competitive or enters the captive-demand regime. For parallel networks, this reduces to the simple summation $\bar{\mu}_{-n, k} = \sum_{m \neq n} \mu_{m, k}$, so no linear program is required. \\ \hline
TLE (Two-Line Element) & A standardized format for satellite orbital parameters published by the U.S. Space Force and updated multiple times daily at celestrak.org. From TLE data, satellite positions and all link capacities are deterministically computable hours to days ahead, enabling proactive regulatory design. \\ \bottomrule
\end{tabular}%
}
\end{table}
\section{Additional Theoretical Results}

This appendix contains the proofs of Theorem~\ref{t2b}, Proposition~\ref{p3}, and Proposition~\ref{p4}, together with the two-operator appendix-only results used in Section \ref{sec4}.

\subsection{Proofs of Theorem~\ref{t2b} and Proposition~\ref{p3}}

\medskip\noindent\textbf{Proof of Theorem~\ref{t2b} (Existence of Nash Equilibrium under a Regulatory Price Cap).}
Suppose the strategy space of each operator is restricted by a finite regulatory price cap, that is, $p_{n}{\in}[0,\bar{p}]$ for all $n=1,{\dots},N$, with $\bar{p}<{\infty}$. If Assumption QC holds, namely that each revenue function ${\pi}_{n}(p_{n},p_{-n})$ is quasi-concave in its own price $p_{n}$ for every fixed $p_{-n}$, then the pricing game admits at least one pure-strategy Nash equilibrium $p^{*}{\in}[0,\bar{p}]^{N}$.

\medskip\noindent\textit{Proof.} Under the finite price cap, the joint strategy space $[0, \bar{p}]^{N}$ is compact and convex. By Theorem~\ref{t1}(c), each operator's revenue function ${\pi}_{n}(p)$ is continuous in the joint price vector $p$. Under Assumption QC, ${\pi}_{n}(p_{n},p_{-n})$ is quasi-concave in $p_{n}$, which implies that the best-response correspondence $BR_{n}(p_{-n})$ is nonempty and convex-valued for every fixed $p_{-n}$. By Berge's Maximum Theorem, the best-response correspondence is upper hemicontinuous. Therefore, the joint best-response correspondence $\prod_{n=1}^{N} BR_{n}(p_{-n})$ satisfies the conditions of Kakutani's fixed-point theorem. Hence, there exists $p^{*}{\in}[0,\bar{p}]^{N}$ such that $p^{*}{\in}\prod_{n=1}^{N} BR_{n}(p_{-n}^{*}),$ which is a pure-strategy Nash equilibrium. $\qed$

\medskip\noindent\textbf{Proof of Proposition~\ref{p3} (Asymptotic Price Divergence).}
Consider an $N$-operator parallel packet routing network in which the aggregate rival capacity satisfies $\bar{{\mu}}_{-n}=D+{\varepsilon},{\varepsilon}{\to}0^{+}.$ Then any interior profit-maximizing equilibrium price of operator $n$ satisfies $p_{n}^{*}({\varepsilon})={\Theta}({\varepsilon}^{-1/2}).$ Moreover, the divergence rate is independent of the number of operators $N$.

\medskip\noindent\textit{Proof.} For the parallel network, the Wardrop equal-cost condition implies that the generalized costs on the used routes are equal. Hence, for operator $n$, $p_{n}^{*}=B-A,$ where $A=\frac{1}{{\mu}_{n}-f_{n}^{*}},B=\frac{1}{\bar{{\mu}}_{-n}-F_{-n}^{*}}.$ The first-order optimality condition for revenue maximization, $\max_{p_{n}} {\pi}_{n}(p_{n})=p_{n}F_{n}^{*},$ yields $p_{n}^{*}=F_{n}^{*}(A^{2}+B^{2}).$ Now consider the asymptotic regime ${\varepsilon}{\to}0^{+}$, where rival capacity approaches the demand threshold from above:

\begin{equation}
\bar{{\mu}}_{-n}=D+{\varepsilon}.
\end{equation}

As the market approaches the regime boundary, the residual rival slack satisfies $\bar{{\mu}}_{-n}-F_{-n}^{*}=O({\varepsilon}^{1/2}),$ so that $B=\frac{1}{\bar{{\mu}}_{-n}-F_{-n}^{*}}={\Theta}({\varepsilon}^{-1/2}).$ By contrast, ${\mu}_{n}-f_{n}^{*}$ remains bounded away from zero in the interior solution, implying $A=O(1).$ Therefore, $A^{2}+B^{2}={\Theta}({\varepsilon}^{-1}),$ and since the equilibrium flow satisfies $F_{n}^{*}={\Theta}({\varepsilon}^{1/2}),$ it follows that $p_{n}^{*}=F_{n}^{*}(A^{2}+B^{2})={\Theta}({\varepsilon}^{1/2}){\Theta}({\varepsilon}^{-1})={\Theta}({\varepsilon}^{-1/2}).$ The argument depends only on the local slack of operator $n$'s rival capacity and therefore does not depend on the number of competing operators $N$. $\qed$

\subsection{Proof of Proposition~\ref{p4} and Corollary~\ref{c3} (N-Operator Equilibrium System)}

\medskip\noindent\textbf{Proof of Proposition~\ref{p4} (Equilibrium System for the $N$-Operator Parallel Network).}
Consider the $N$-operator parallel packet routing network in Regime A. Define the relative flow ratios $r_{n}=\frac{f_{n}}{f_{1}},n=1,{\dots},N,$ with normalization $r_{1}\ensuremath{\equiv}1$. Let $S(r)=\sum_{m=1}^{N} r_{m},f_{n}(r)=\frac{D r_{n}}{S(r)},$ and define $A_{n}(r)=\frac{1}{{\mu}_{n}-f_{n}(r)}.$ For each operator $n$, let $S_{-n}(r)=\sum_{m{\ne}n}^{} A_{m}(r)^{-2},$ and $C_{n}(r)=A_{n}(r)^{2}+\frac{1}{S_{-n}(r)}.$ Then the first-order condition for operator $n$ is $p_{n}^{*}=f_{n}(r) C_{n}(r).$ Combining the Wardrop equal-cost condition with the first-order pricing conditions yields the following $\left(N - 1\right)$-dimensional nonlinear equilibrium system:
\begin{equation}
    (t_n - t_1) + A_n (r) - A_1 (r) = f_1 (r) C_1 (r) - f_n (r) C_n (r), \quad n = 2, \dots, N.
    \label{eq8}
\end{equation}

The following statements hold:

(i) Equation \eqref{eq8} admits a unique solution $r^{*}$ in Regime A.

(ii) When $N=2$, Equation \eqref{eq8} reduces to the scalar equilibrium equation presented in Appendix C.

(iii) When $N{\ge}3$, no further closed-form reduction is generally available, and the equilibrium is obtained numerically using Algorithm \ref{alg3}.

\begin{corollary}[Uniqueness of the $N$-Operator Equilibrium Flow Ratios]\label{c3}
For any $N$-operator parallel network in Regime A satisfying $\sum_{n=1}^{N} {\mu}_{n}>D,$ the equilibrium system given by Equation \eqref{eq8} admits a unique solution $r^{*}{\in}\mathbb{R}_{>0}^{N-1}.$
\end{corollary}

\medskip\noindent\textit{Proof.} By Theorem~\ref{t1}(a), the Wardrop equilibrium flow vector $f^{*}(p^{*};{\mu}_{k})$ is unique for any equilibrium price vector $p^{*}$. Since the flow ratios are defined by $r_{n}^{*}=\frac{f_{n}^{*}}{f_{1}^{*}},n=2,{\dots},N,$ the corresponding ratio vector $r^{*}$ is uniquely determined by the equilibrium flow vector. Therefore, Equation \eqref{eq8} admits at most one solution. To establish existence, Theorem~\ref{t1}(d) guarantees the existence of a Nash equilibrium price vector $p^{*}$ under Assumption QC. Substituting the corresponding unique Wardrop equilibrium flow $f^{*}(p^{*};{\mu}_{k})$ into the definition of $r_{n}$ yields a feasible ratio vector $r^{*}$ satisfying Equation \eqref{eq8}. Hence, a unique solution exists.

For $N=2$, Equation \eqref{eq8} reduces to a one-dimensional system because $S_{-1}=A_{2}^{-2},S_{-2}=A_{1}^{-2},$ which implies $C_{1}=C_{2}=A_{1}^{2}+A_{2}^{2}.$ Hence, Equation \eqref{eq8} collapses to the scalar equilibrium equation given in Appendix C. For $N{\ge}3$, the terms $S_{-n}(r)$ depend on all relative flow ratios through $S(r)$, which couples the system across all operators. As a result, no further closed-form reduction is generally available, and numerical solution is required. $\ensuremath{\square}$

\section{Two-Operator Closed-Form Results}

This appendix derives the closed-form results for the $N = 2$ special case referenced in Section~\ref{sec4}. These results are the satellite routing analogues of the explicit toll formulas for the two-route road network \cite{Pigou1920, yang2004multi}.

\subsection{Scalar Equilibrium Equation and Cutoff Price}

For the two-operator case $\left(N = 2\right)$, the equilibrium system in Equation \eqref{eq8} reduces to a one-dimensional nonlinear equation in the relative flow ratio $r=\frac{f_{2}}{f_{1}}.$ Let $A(r)=\frac{1}{{\mu}_{1}-f_{1}(r)},B(r)=\frac{1}{{\mu}_{2}-f_{2}(r)},$ where $f_{1}(r)=\frac{D}{1+r},f_{2}(r)=\frac{Dr}{1+r}.$ Since $N=2$, the curvature terms satisfy $C_{1}(r)=C_{2}(r)=A(r)^{2}+B(r)^{2}.$ Substituting these expressions into Equation \eqref{eq8} yields the scalar equilibrium equation
\begin{equation}
    (t_2 - t_1) + B(r) - A(r) = \frac{D(1-r)}{1+r} (A(r)^{2} + B(r)^{2}).
    \label{eq9}
\end{equation}
By Corollary~\ref{c3}, Equation \eqref{eq9} admits a unique solution $r^{*}>0$. Once $r^{*}$ is obtained, the Nash equilibrium flows are recovered as $f_{1}^{*}=\frac{D}{1+r^{*}},f_{2}^{*}=\frac{Dr^{*}}{1+r^{*}},$ and the corresponding equilibrium prices are $p_{n}^{*}=f_{n}^{*}(A^{*2}+B^{*2}),n=1,2.$ The unique root $r^{*}$ can be computed in $O(\log\frac{1}{{\varepsilon}})$ time using Brent's method, as implemented in Algorithm \ref{alg5}.

\begin{proposition}[Cutoff Price in the Two-Operator Case]\label{p1}
Suppose ${\mu}_{1},{\mu}_{2}>D$. The cutoff price at which operator 1's contested flow becomes zero is:
\begin{equation}
\bar{p}_{1}(p_{2})= (t_2 - t_1) + \frac{1}{{\mu}_{2}-D} - \frac{1}{{\mu}_{1}}.
\end{equation}
At any price $p_{1}{\ge}\bar{p}_{1}(p_{2}),$ operator 1 receives zero contested flow and therefore earns zero revenue from competitive demand. In Regime B, however, operator 1 retains captive demand satisfying $f_{1}^{*}{\ge}{\chi}_{1}>0,$ so this cutoff boundary no longer applies.
\end{proposition}

\subsection{Sufficient Quasi-Concavity Condition and Regulatory Price Cap}

For the two-operator case, consider operator 1's revenue function ${\pi}_{1}(p_{1},p_{2})=p_{1} f_{1}^{*}(p_{1},p_{2}).$ Differentiating twice with respect to $p_{1}$ yields:
\begin{equation}
\frac{{\partial}^{2} {\pi}_{1}}{{\partial} p_{1}^{2}}= \frac{2}{(A^{2}+B^{2})^{2}} \left[ p_1 (B^{3}-A^{3}) - (A^{2}+B^{2}) \right].
\end{equation}
A sufficient condition for quasi-concavity is therefore:
\begin{equation}
    \Upsilon_1 := \frac{p_1 (B^3 - A^3)}{A^2 + B^2} \le 1.
\end{equation}

To obtain a closed-form bound, we evaluate the worst-case configuration as $f_{1}{\to}0$, under which

\begin{equation}
A{\to}A_{0}=\frac{1}{{\mu}_{1}},B{\to}B_{0}=\frac{1}{{\mu}_{2}-D}.
\end{equation}

This yields the following sufficient condition.

\begin{theorem}[QC-SC: Sufficient Condition for Quasi-Concavity in the Two-Operator Case]\label{thm:qcsc}

In Regime A, define $A_{0}=\frac{1}{{\mu}_{1}},B_{0}=\frac{1}{{\mu}_{2}-D},C_{0}=\frac{1}{{\mu}_{2}},D_{0}=\frac{1}{{\mu}_{1}-D},$ and let ${\Delta}t=t_{2}-t_{1}.$ For operator 1, define $K_{1}=\begin{cases}\frac{B_{0}^{3}-A_{0}^{3}}{A_{0}^{2}+B_{0}^{2}}, & B_{0}>A_{0}, \\ 0, & \text{otherwise},\end{cases}$ and define $K_{2}$ analogously for operator 2. If the regulatory price cap satisfies
\begin{equation}
    \bar{p} \le \bar{p}^{\mathrm{QC}}, \bar{p}^{\mathrm{QC}} = \min \{\bar{p}^{\mathrm{QC_1}}, \bar{p}^{\mathrm{QC_2}}\}
\end{equation}
where
\begin{align}
    \bar{p}^{\mathrm{QC_1}} &= \frac{1}{K_1} - (B_0 - A_0) - \Delta t \\
    \bar{p}^{\mathrm{QC_2}} &= \frac{1}{K_2} - (D_0 - C_0) + \Delta t
\end{align}
then Assumption QC holds simultaneously for both operators over the strategy space $[0, \bar{p}]^{2}$. Moreover, the bound $\bar{p}^{\mathrm{QC}}$ is computable in $O(1)$ time for each orbital window from the TLE-derived capacity parameters.
\end{theorem}

\begin{corollary}[Double Squeeze Effect]\label{c2}
Let ${\sigma}_{2}={\mu}_{2}-D.$ As ${\sigma}_{2}{\to}0^{+},$ the QC-safe regulatory cap satisfies $\bar{p}^{\mathrm{QC}}{\to}0,$ while the equilibrium price diverges according to $p_{1}^{*}\ensuremath{\sim}{\sigma}_{2}^{-1/2}{\to}{\infty}.$ Hence, the admissible price cap shrinks to zero precisely when the endogenous market incentive to raise prices becomes unbounded.

For the two-operator case, the explicit expression for the QC-safe cap is
\begin{equation}
    \bar{p}^{\mathrm{QC}} = \frac{\mu_1 \sigma (\sigma^2 + \mu_1^2)}{(\mu_1 - \sigma)(\sigma^2 + \mu_1\sigma + \mu_1^2)} - \frac{\mu_1 - \sigma}{\mu_1 \sigma} - \Delta t, \sigma = \mu_2 - D.
\label{eq16}
\end{equation}
Moreover, the cap is monotone in the rival slack parameter:

\begin{equation}
\frac{{\partial}\bar{p}^{\mathrm{QC}}}{{\partial}{\sigma}}>0.
\end{equation}

Therefore, as the rival capacity margin tightens, the allowable QC-safe cap decreases monotonically.
\end{corollary}

The welfare-optimal regulatory cap is obtained by solving

\begin{equation}
\min_{\bar{p}{\ge}0}TC(\bar{p})
\end{equation}

subject to

\begin{equation}
\bar{p}{\le}\bar{p}^{\mathrm{QC}}({\mu}_{k}),
\end{equation}

which is computed using Algorithm \ref{alg4}.

\bibliographystyle{elsarticle-harv}
\bibliography{cas-refs}

\end{document}
